% !TeX encoding = UTF-8
% !TeX program = pdflatex
% !BIB program = bibtex

%%% To write an article in English, please use the option ``english'' in order
%%% to get the correct hyphenation patterns and terms.
%%% \documentclass[english]{class}
%%% for anonymizing an article you can use the ``anonymous'' option.
%%%
%%% Um einen Artikel auf deutsch zu schreiben, genügt es die Klasse ohne
%%% Parameter zu laden.
%%% Zur Anonymisierung kann die ``anonymous'' Option genutzt werden.
\documentclass[english]{lni}
%%
\usepackage{wasysym}
\usepackage{graphicx}
\usepackage{balance}  % for  \ance command ON LAST PAGE  (only there!)
\usepackage{amsmath}
%\usepackage{amssymb}
\usepackage{microtype}
\usepackage{booktabs}
\usepackage{wrapfig}
\usepackage{enumitem}
\usepackage{listings}
\usepackage[english]{babel}
\usepackage{multirow}
\usepackage{times}
\usepackage{tabularx}
\usepackage{verbatim}
\usepackage{caption}
\usepackage{subcaption}
\usepackage[export]{adjustbox}
\usepackage[utf8]{inputenc}
\usepackage[ruled,vlined,linesnumbered]{algorithm2e}%
\usepackage{etoolbox}
\usepackage{tcolorbox}
\usepackage{xcolor}
\usepackage[colorinlistoftodos]{todonotes}
\usepackage{bbm}
\usepackage[flushleft]{threeparttable}
\usepackage{wrapfig}
\usepackage{ulem}
\DeclareMathOperator{\poly}{\textup{poly}}


\newcommand{\code}[1]{{\textsf{{\small #1}}}}
% \newcommand{\todo}[1]{\textcolor{red}{\textit{#1}}}
\newcommand{\todomargin}[1]{\marginpar{\textcolor{red}{\textit{#1}}}}

\definecolor{forestgreen}{HTML}{228C22}
\definecolor{myblue}{HTML}{1E88E5}
\definecolor{myyellow}{HTML}{004D40}
\definecolor{mymagenta}{HTML}{D81B60}

\renewcommand{\labelitemi}{$\circ$}
\renewcommand{\labelitemii}{$-$}

\def\Snospace~{\S{}}
\newcommand\sectionautorefname{\Snospace}
\newcommand\subsectionautorefname{\Snospace}
\renewcommand*\algorithmautorefname{Alg.}
\newcommand*\exampleautorefname{Example}
\newcommand*\lemmaautorefname{Lemma}
%\newcommand*\theoremautorefname{Theorem}
\newcommand*\propertyautorefname{Property}
\newcommand*\problemautorefname{Problem}
\newcommand*\definitionautorefname{Def.}
\renewcommand*{\algorithmcflinename}{line}
% \renewcommand\equationautorefname{Eq.}
% \renewcommand*\figureautorefname{Fig.}

\makeatletter
\newcommand{\removelatexerror}{\let\@latex@error\@gobble}
\makeatother

\SetAlgoSkip{}

\newcommand{\exampleboxh}[2]{
\begin{tcolorbox}[colback=white,
	float,
	floatplacement=h,
	enlarge top by=-.6em,
	enlarge bottom by=-.6em,
	colframe=gray!75!black,
	title=\scshape\fontseries{sb}\selectfont#1,
	boxsep=3pt,
	left=1pt,
	right=1pt,
	top=0pt,
	bottom=0pt,
	toptitle=0pt,
	bottomtitle=0pt,
	lefttitle=1pt,
	righttitle=1pt,
	]{\small%
#2%
}
\end{tcolorbox}%
}

\newcommand{\examplebox}[2]{
\begin{tcolorbox}[colback=white,
	float,
	floatplacement=t,
	enlarge bottom by=-.6em,
	colframe=gray!75!black,
	title=\scshape\fontseries{sb}\selectfont#1,
	boxsep=3pt,
	left=1pt,
	right=1pt,
	top=1pt,
	bottom=1pt,
	toptitle=0pt,
	bottomtitle=0pt,
	lefttitle=1pt,
	righttitle=1pt,
	]{\small%
		#2%
	}
\end{tcolorbox}%
}

\newtheorem{corollary}{Corollary}
\newtheorem{problem}{Problem}
\newtheorem{lemma}{Lemma}
\newtheorem{definition}{Definition}
\newtheorem{example}{Example}
\newtheorem{property}{Property}
\newtheorem{theorem}{Theorem}

\newcommand{\qed}{\leavevmode
  \hbox to.77778em{%
  \hfil\vrule
  \vbox to.675em{\hrule width.6em\vfil\hrule}%
  \vrule\hfil}}

  \lstset{ %
  belowskip=-2em,
  backgroundcolor=\color{white},   % choose the background color; you must
  basicstyle=\footnotesize\ttfamily,        % the size of the fonts that are
  %used
  %for the code
  breakatwhitespace=false,         % sets if automatic breaks should only
  %happen at whitespace
  breaklines=true,                 % sets automatic line breaking
  captionpos=b,                    % sets the caption-position to bottom
  commentstyle=\color{Brown},    % comment style
  deletekeywords={...},            % if you want to delete keywords from the
  %given language
  escapeinside={\%*}{*)},          % if you want to add LaTeX within your code
  extendedchars=true,              % lets you use non-ASCII characters; for
  %8-bits encodings only, does not work with UTF-8
  frame=,                    % adds a frame around the code
  keepspaces=true,                 % keeps spaces in text, useful for keeping
  %indentation of code (possibly needs columns=flexible)
  keywordstyle=\color{ACMDarkBlue},       % keyword style
  language=,                 % the language of the code
  morekeywords={*,RETURN,SEQ,WITHIN,PATTERN,NOT,
      WHERE,AND,NSEQ,ms,h, min, range},
  %
  %if you want to add more keywords to the set
  numbers=none,                    % where to put the line-numbers; possible
  %values are (none, left, right)
  numbersep=5pt,                   % how far the line-numbers are from the
  %code
  numberstyle=\tiny\color{gray}, % the style that is used for the line-numbers
  rulecolor=\color{black},         % if not set, the frame-color may be
  %changed
  %on line-breaks within not-black text (e.g. comments (green here))
  showspaces=false,                % show spaces everywhere adding particular
  %underscores; it overrides 'showstringspaces'
  showstringspaces=false,          % underline spaces within strings only
  showtabs=false,                  % show tabs within strings adding
  %particular
  %underscores
  stepnumber=1,                    % the step between two line-numbers. If
  %it's
  %1, each line will be numbered
  stringstyle=\color{RedOrange},     % string literal style
  tabsize=2,                       % sets default tabsize to 2 spaces
  title=\lstname,                   % show the filename of files included
  %with
  %\lstinputlisting; also try caption instead of title
}


% \hypersetup{
%   pdftoolbar=true,        % show Acrobat?s toolbar?
%   pdfmenubar=true,        % show Acrobat?s menu?
%   pdffitwindow=false,     % window fit to page when opened
%   pdfstartview={FitH},    % fits the width of the page to the window
%   pdftitle={},    % title
%   pdfauthor={},     % author
%   pdfsubject={},   % subject of the document
%   pdfcreator={},   % creator of the document
%   pdfproducer={}, % producer of the document
%   pdfkeywords={}, % list of
%   %keywords
%   pdfnewwindow=true,      % links in new window
%   colorlinks=true,       % false: boxed links; true: colored links
%     linkcolor=Brown,          % color of internal links (change box color
% %	with
% %	linkbordercolor)
%     citecolor=OliveGreen,        % color of links to bibliography
%     filecolor=magenta,      % color of file links
%     urlcolor=NavyBlue           % color of external links
% }

\newcommand{\sys}{DISCES\xspace}

\renewcommand{\UrlFont}{\ttfamily\scriptsize} %?
\newcommand{\sstitle}[1]{\smallskip\noindent\textbf{#1.\/}}
\newcommand{\sstitlenoskip}[1]{\noindent\textbf{#1.\/}}

\newcommand{\change}[1]{\textcolor{blue}{#1}}


\DeclareMathOperator{\dom}{dom}
\DeclareMathOperator{\argmin}{argmin}

\newcommand\Mark[1]{\textsuperscript#1}
\newcommand{\True}{\mathit{True}}
\newcommand{\False}{\mathit{False}}
\def\ua{\Mark{*}}
\def\ub{\Mark{\dag}}
\def\uc{\Mark{\#}}
\def\ud{\Mark{+}}


\begin{document}

\def\Snospace~{\S{}}
\renewcommand\sectionautorefname{\Snospace}
\renewcommand\subsectionautorefname{\Snospace}


%%% Mehrere Autoren werden durch \and voneinander getrennt.
%%% Die Fußnote enthält die Adresse sowie eine E-Mail-Adresse.
%%% Das optionale Argument (sofern angegeben) wird für die Kopfzeile verwendet.
\title[Incremental Updates of Discovered Event Queries]{Embracing Change:
Incremental Updates of Discovered Event Queries}
%% \subtitle{Untertitel / Subtitle} % if needed
 \author[1]{Rebecca Sattler}{sattlere@hu-berlin.de}{0000-0002-7342-7131}
 \author[1]{Sarah Kleest-Meißner}{kleest-meissner.sarah@alumni.hu-berlin.de}{0000-0002-4133-7975}
 \author[1]{Steven Lange}{langestx@hu-berlin.de}{0009-0001-7217-8188}
 \author[1]{Markus Schmid}{mlschmid@mlschmid.de}{0000-0001-5137-1504}%
 \author[1]{Nicole
 Schweikardt}{schweikn@informatik.hu-berlin.de}{0000-0001-5705-1675}%
 \author[1]{Matthias
 Weidlich}{matthias.weidlich@hu-berlin.de}{0000-0003-3325-7227}%

 \affil[1]{Humboldt-Universität zu Berlin\\Institut für Informatik\\Unter
 den Linden 6\\10099 Berlin\\Germany}

 \maketitle
\begin{abstract}
In complex event processing (CEP), queries are evaluated continuously over
streams of events to detect situations of interest, thereby
facilitating reactive applications.
However, users often lack insights into the precise event pattern that
characterizes the situation, which renders the
definition of the respective queries challenging.
Once a database of finite, historic
streams, each containing a materialization of the situation of interest, is
available, query discovery supports users in the
definition of the desired queries. It constructs the queries that match a
certain share of the given streams, as determined by a support
threshold. Yet, upon changes in the database or changes of the support
threshold, existing algorithms need to construct the resulting
queries from scratch, neglecting the queries obtained in previous runs.
In this paper, we aim to avoid the resulting inefficiencies by techniques
for incremental query discovery. We first provide a theoretical analysis of
the problem context, before presenting algorithmic solutions to cope with
changes in the stream database or the adopted support threshold.
Our experiments using real-world data show that our incremental query
discovery
reduces the runtimes by up to three orders of magnitude compared to a
baseline solution.
\end{abstract}
\begin{keywords}
Complex Event Processing \and Event Streams \and Query Discovery \and
Incremental Updates
\end{keywords}
%%% Beginn des Artikeltexts

\section{Introduction}
\label{sec:introduction}

%\todo[inline]{R2 on queries for outliers: The stream
%database is assumed to represent
%deviating/abnormal/outlier behavior,
%so that the discovered queries capture the common
%patterns that are inherent to the respective
%situations.
%We intend to clarify this in the paper.}
Complex event processing (CEP) is a computational paradigm
to facilitate reactive and proactive
applications~\cite{DBLP:journals/csur/CugolaM12,DBLP:journals/vldb/GiatrakosAADG20}.
 In CEP, a set of queries are evaluated continuously over
a stream of events in order to detect situations of
interest in domains, such as urban
transportation~\cite{DBLP:conf/edbt/ArtikisWSBLPBMKMGMGK14},
computational
finance~\cite{DBLP:journals/computer/ChandramouliAGSR10},
and supply chain
management~\cite{DBLP:journals/cii/KonovalenkoL19}.
Often, these situations of interest are deviating or
abnormal behaviour, so that their timely detection
facilitates the implementation of counter-measures.

In essence, CEP queries capture sequential patterns of
events, i.e., they describe a sequence of events of
specific types, additional conditions on the attribute
values of the respective events, and a time window for
their occurrence. As the exact characterization of the
event patterns is challenging in practice, one may support
users in the specification of queries
with techniques for automated query
discovery~\cite{icep,ilminer,DBLP:conf/btw/Kleest-Meissner23}.
 Here, the assumption is
that a database of finite, historic (sub-)streams that
cover the situation of interest is available. That is,
if the situation of interest is some deviating or
abnormal behaviour, the assumption is that the
database contains only the respective streams, not the
common behaviour. Then,
queries that match a certain share of these streams, as
determined by a support threshold, can be extracted and
subsequently be assessed and validated by a user.

\begin{figure}[t]
	\centering
	\includegraphics[width=0.8\columnwidth]{img/overview.drawio.pdf}
	%\vspace{-1.0em}
	\caption{The general setting of incremental event
	query discovery: An update of the stream
	database or support threshold shall be incorporated in
	an updated set of queries.}
	\label{fig:overview}
	\vspace{-1em}
\end{figure}

However, over time, the database of streams used as a
basis for query discovery may be subject to change.
Additional streams that capture a situation of interest
may become available, or streams that are already in the
database may be considered outdated. As part of that, the
requirements for query discovery in terms of the adopted
support threshold may also change, i.e., it may be relaxed
to increase the number of discovered queries, or
strengthened to avoid false positives in the result.

Yet, upon changes in the database or in the support
threshold, existing discovery algorithms need to construct
the resulting queries from scratch. This is problematic
since solving the problem of event query discovery is
computationally hard, as an exponentially growing space of
candidate queries (in the query length and attribute
domains) needs to be explored and the evaluation of a
single query also shows an exponential runtime complexity
(in the number of variables, i.e., conditions over
multiple events)~\cite{icdt2022}.

In this paper, we set out to avoid
inefficient recomputation with strategies for incremental
event query discovery. Our idea is that the discovery
result shall be reused and adapted once the databases or
the support threshold change, as illustrated in
\autoref{fig:overview}. To this end, after giving
formal preliminaries (\autoref{sec:problem-setting}), we
make the following contributions:
\begin{enumerate}[left=.2em,nosep,label=(\roman*)]
\item We present a theoretical analysis of
the problem context (\autoref{sec:updates}).
Specifically, we show that under certain conditions, the
result set of queries for an updated stream database or an
updated support threshold denotes a subset of the existing
result.
\item We present algorithmic
solutions for incremental query discovery
(\autoref{sec:algos}). They adapt a given set of queries,
thereby aiming to avoid inefficient recomputation.
\end{enumerate}
We assess the feasibility of our
solution in several experiments
(\autoref{sec:evaluation}), showing that incremental
discovery reduces the runtimes by up to 10x.
Finally, we review related work
(\autoref{sec:related-work}) and conclude the paper
(\autoref{sec:conclusions}).








\section{Background on Query Discovery}
\label{sec:problem-setting}
%\todo[inline]{R2: The explanation of the query model is good. However, it
%could be supplemented by different pattern-matching semantics
%(contiguous, skip-till-next, skip-till-any) and an explanation of how the
%window parameter can be learned.
%
%A: We will add a statement on how to incorporate further query semantics and
%how to determine the window size.}
Below, we introduce preliminary notions to study
the problem of event query discovery.

%\subsection{Event Streams and Queries}
%\label{sec:streams}

%We adopt some standard notation for sequences, as follows. For a set
%${Y}$, the set of all sequences over ${Y}$ is ${Y}^*$.
%For a finite sequence $z=\langle y_1,\ldots, y_n \rangle \in {Y}^*$, we
%write $z[i] = y_i$ for the $i$-th element.

\sstitle{Event streams}
We based our work on the notion of a multivariate event
stream. Such a stream has an \emph{event schema}, captured
by a tuple of attributes $\mathcal{A}= (A_1, \ldots,
A_n)$, where $\dom(A_i)$ is the finite domain of attribute
$A_i$, $1\leq i\leq n$. All events have the same schema
and we assume that attribute domains are disjoint,
$\cap_{i=1}^n \dom(A_i) =
\emptyset$. As a short-hand, we write $\Gamma =
\bigcup_{i=1}^n{\dom(A_i)}$ for the alphabet defined by
all possible values of attributes. Moreover, an
\emph{event} is an instantiation of the event schema,
captured as a tuple $e=(a_1, \ldots, a_n)$ with $a_i\in
\dom(A_i)$ for $1\leq
i\leq n$. By $e.A_i$, we denote the value of attribute
$A_i$ of event $e$.

An \emph{event stream} is a finite
sequence of events $S=\langle e_1, \ldots, e_l
\rangle$, i.e., a finite subsequence of a
potentially unbounded sequence of events. We write $|S|$
for the length of the stream, and $S[i]$ for its $i$-th
event. A \emph{stream database} is a finite set of
streams $D=\{S_1,\ldots, S_d\}$.

\sstitle{Event queries}
We adopt a query model that captures sequences of events,
which are correlated by predicates over their attribute
values~\cite{DBLP:conf/sigmod/ZhangDI14}. We neglect the
definition of an explicit time window, as it may be
derived from the maximal time difference between events in
one of the streams. That is, taking the time between the start and end of
each stream, the maximal time span over all streams is an upper bound for the
time interval, in which a situation of interest is expected to occur.
Moreover, we
follow~\cite{DBLP:conf/btw/Kleest-Meissner23} and use a
linearized representation of event queries.
Let $\mathcal{A}= (A_1, \ldots,
A_n)$ be an event schema, $\mathcal{X}$ be a finite set of
variables
such that $\mathcal{X}\cap
\Gamma=\emptyset$, and $\_ \notin \mathcal{X}\cup
\Gamma$ be a placeholder symbol. Then, we define a
\emph{query
term} $t=(u_{1},\ldots,u_{n})$ as an $n$-tuple comprising
attribute
values, variables, and placeholders:
$u_{j}\in \dom(A_j) \cup \mathcal{X} \cup \{ \_ \}$,
$1\leq j\leq n$, while we prohibit terms containing only
placeholders, i.e., $u_{j}\in
\dom(A_j) \cup \mathcal{X}$ for at least one $1\leq j\leq n$.
By $\varepsilon$, we denote the empty query term that
contains only
placeholders, i.e., %$\varepsilon=(\underbrace{\_,\dots,\_}_{n})$
$\varepsilon=(\_,\dots,\_)$.
Based thereon, an \emph{event
query} is a finite sequence of query terms, $q=\langle
t_1,\ldots,t_k \rangle$,
with $q[i]$ denoting the $i$-th term. In a query, variables need to occur in at
least two query terms for the same attribute, i.e., for any term
$q[i]=(u_{1},\ldots,u_{n})$ with $u_j\in \mathcal{X}$,
$1\leq i\leq k$ and $1\leq j\leq n$, there exists another query term
$q[p]=(u'_{1},\ldots,u'_{n})$ with $u_j = u'_j$,
$1\leq p\leq k$ and $p\neq i$. The empty query is
$\langle\varepsilon\rangle$, i.e., the query containing
the empty query term. By $\mathcal{Q}$, we denote the
universe of all queries.

From an intuitive point of view, a query term describes a
single event that shall be matched by defining a specific
value, variable, or placeholder per attribute. While a
placeholder denotes the absence of any constraint, the use
of variables captures the correlation of events as
matching events need to have equivalent values for the
respective attributes.


%\item $\tau\in \mathbb{N}$ is a window with $\tau \leq k$.
%\end{itemize}
%Given an event
%schema $\mathcal{A} =
%(A_1, \ldots, A_n)$ and a finite set of variables $\mathcal{X}$,
%an event query
%%is a tuple $q=(T,\tau)$, where
%%\begin{itemize}[nosep]
%%\item
%is a finite sequence of query terms $q=\langle t_1,\ldots,t_k \rangle$,
%with $q[i]$ denoting the $i$-th term. Each term
%$t_i=(u_{1},\ldots,u_{n})$, $1\leq i\leq k$, is an $n$-tuple built of
%attribute
%values, variables, and the dedicated
%symbol $\varepsilon \notin \Gamma \cup \mathcal{X}$ to capture the absence of
%a
%constraint,
%i.e.,
%$u_{j}\in \dom(A_j) \cup
%\mathcal{X} \cup \{\varepsilon\}$, $1\leq j\leq n$.
%%\item $\tau\in \mathbb{N}$ is a window with $\tau \leq k$.
%%\end{itemize}

%Note that variables that occur only once in a query are
%semantically
%equivalent
%to placeholders and, as such, will be largely neglected in the remainder.
%Also,
%only the elements of $\mathcal{X}$ that occur at least twice in a query are
%referred to as variables.
%Adopting this terminology,

Moreover, we distinguish three classes of queries: A
\emph{type query} contains only query terms defining
attribute values (and
placeholders). A \emph{pattern query} contains only query
terms defining variables (and
placeholders). All other queries are referred to as
\emph{mixed queries}.

% only
%variables, or a combination of both.
%These are denoted as type queries, pattern queries, and mixed queries,
%respectively.
% To improve readability, the non specified attributes are left empty in the query
%terms.

\begin{table}[t]
	\caption{Overview of notations.}
	\label{tab:notations}
	\vspace{-.5em}
	\footnotesize
	\begin{tabularx} {\textwidth}{l
	 @{\hspace{.5em}} p{10cm}}
	\toprule
	Notation & Explanation \\
	\midrule
	$\mathcal{A}=(A_1, \ldots, A_n )$ & An event
	schema, a
	tuple of attributes\\
	$\Gamma  = \bigcup_{i=1}^n{\dom(A_i)}$ & The stream alphabet\\
	$e= (a_1, \ldots, a_n)$ & An event, a tuple
	of attribute values\\
	$e.A_i$ & The value $a_i$ of attribute $A_i$ of
	event $e$\\
	%$e.t$ & The occurrence time of event $E$\\
	$S=\langle e_1, \ldots, e_l \rangle$ & An event
	stream, a finite sequence of events\\
	$S[i]$ & The $i$-th event of the event stream $S$\\
	$D=\{S_1,\ldots, S_d \}$ & A stream database, a
	set of event streams\\
%	$\Gamma_D$ & The supported stream alphabet of stream database $D$\\
	\midrule
	$\mathcal{X}$ & A finite set of query variables, $\mathcal{X} \cap
	\Gamma=\emptyset$\\
	$\_$ & A placeholder symbol, $\_ \notin \mathcal{X}\cup
	\Gamma$\\
	$t=(u_{1},\ldots,u_{n})$& A query term, an
	$n$-tuple of attribute values, variables, and placeholders\\
%	with $u_{j}\in
%	\dom(A_j) \cup \mathcal{X} \cup \{\_ \}$, $1\leq j\leq n$, such that not
%	all its components are placeholders\\
	$\varepsilon=(\_,\ldots,\_)$& The empty query term, an
$n$-tuple of placeholders\\
	$q=\langle
	t_1, \ldots, t_k \rangle$ & An event query, a finite sequence of
	query terms \\
	$q[i]$ & The $i$-th term of the query $q$\\
	$\langle \varepsilon\rangle$ & The empty query containing only the empty
	query term \\
	\midrule
	$S \models q$ & The stream $S$ supports the
	query $q$, i.e., there exists a match\\
	$M_+(q,D)$ & The set of streams in $D$ that match the query $q$\\
%	$M_-(q,D)$ & The set of streams in $D$ that do not match the query $q$\\
	supp$(q,D)$ & The support of query $q$ in stream database $D$\\
	$\theta$ & The support threshold, a real number in the interval $(0,1]$\\
	$s(|D|,\theta)$ & The stream threshold, the min number of streams
	to match to fulfill the support threshold\\
	$\Gamma(D,\theta)$ & The frequent stream alphabet\\
	$Q_{D}^{\theta}$ & The set of all matching queries\\
	%$q$ in stream $S$\\
	\bottomrule
	\end{tabularx}
	\vspace{-1em}
\end{table}

We define the semantics of a query $q=\langle
t_1,\ldots,t_k \rangle$ over a stream $S=\langle e_1,
\ldots , e_l\rangle$ as follows: A \emph{match} of $q$ in
$S$ is an injective mapping
$m: \{1,\ldots,k\}\rightarrow \{1, \ldots, l\}$,
such that:
\begin{itemize}[nosep,left=.2em .. 1.5em]
	\item for each query term $q[i]=(u_{1},\ldots,u_{n})$, $1\leq i\leq k$, the
	mapped event $S[m(i)]=(a_1,\ldots, a_n)$ has the
	required attribute
	values, $u_j = a_j$ or $u_j\in \mathcal{X}\cup \{\_ \}$
	for $1\leq 	j\leq n$;
	\item variables are bound to the same attribute values, i.e., for query
	terms $q[i]=(u_1,\ldots, u_n)$ and $q[p]=(u'_1,\ldots, u'_n)$, $1\leq
	i,p\leq k$, it holds that $u_j\in \mathcal{X}$, $1\leq j\leq n$, with
	$u_j=u'_j$ implies that $S[m(i)].A_j = S[m(p)].A_j$; and
 	\item the order of events in the stream is preserved,
 	i.e., $m(i)<m(p)$ for $1\leq
	i<p\leq k$.
\end{itemize}
Note that the set of matches as induced by the above definition corresponds
to a greedy event selection
policy (also known as unconstrained or skip-till-any-match). It does not
limit how often an event can participate in matches. However, in the
remainder, we do not explicitly refer to the set of matches, but consider
only the existence of a match per stream. Specifically, we say that
$q$ matches $S$ and write $S \models q$, if there exists a
match for query $q$ in stream $S$. \autoref{tab:notations}
provides an overview of the introduced notations.

\sstitle{Support} To reason about the relation of a query
$q$ and a stream database $D$, we
define $M_+(q,D)$ as the
set of matching streams, i.e., $M_+(q,D) = \{S \in D \mid
S \models q\}$.
The \emph{support} supp$(q,D)$ of a query is then given as
the fraction of streams in $D$ that match the query,
supp$(q,D)= \frac{|M_+(q,D)|}{|D|}$.
A \emph{support threshold} $\theta$ is a real number in
the interval $(0,1]$.
A query $q$ is \emph{frequent} in $D$, if supp$(q,D) \geq
\theta$.
From the aforementioned notions, we derive the
\emph{stream threshold} $s(|D|,\theta)$ as the minimum
number of streams
that must match to a query in order for the query to be
considered frequent, i.e., $s(|D|,\theta) = \lceil |D| \cdot
\theta \rceil$.
Based thereon, we also derive the frequent stream alphabet, defined by all
attribute values that appear in at least the number of streams as
determined by the stream threshold, i.e., $\Gamma(D,\theta) = \{a \in
\Gamma \mid |\{S \in D \mid \exists \ j\in\{1,\dots,|S|\},
i\in\{1,\dots,n\}: S[j].A_i = a\}| \geq s(|D|,\theta)\}$.
Furthermore, as a short-hand,
we define the set of all matching queries as $Q_D^{\theta}
= \{q \mid \text{supp}(q,D) \geq \theta\}$.

For a given stream database, it is clear that support
thresholds induce a containment relation for the sets of
matching queries. As we later rely on this property, we
state it as follows:

\begin{property}
\label{prop:subset}
Let $D$ be a stream database and $\theta_1$, $\theta_2$ be support thresholds. %with $\theta_1 \leq \theta_2$.
Let $Q_D^{\theta_1}$ and $Q_D^{\theta_2}$ be the sets of all matching queries. %, where $Q_i = \{q \mid \text{supp}(q,D) \geq \theta_i\}$ for $i \in \{1,2\}$.
If $\theta_2 \geq \theta_1$, then $Q_D^{\theta_2} \subseteq Q_D^{\theta_1}$.
\end{property}
% \textbf{Proposition:} Let $D$ be a stream database and $\theta_1$, $\theta_2$ be support thresholds. %with $\theta_1 \leq \theta_2$.
% Let $P_D^{\theta_1}$ and $P_D^{\theta_2}$ be the sets of all matching queries. %, where $Q_i = \{q \mid \text{supp}(q,D) \geq \theta_i\}$ for $i \in \{1,2\}$.
% If $\theta_2 \geq \theta_1$, then $P_D^{\theta_2} \subseteq P_D^{\theta_1}$. \\
\textit{Proof.}
    Let $q \in Q_D^{\theta_2}$. Then, supp$(q,D) \geq
    \theta_2 \geq \theta_1$, and thus, $q \in
    Q_D^{\theta_1}$
    and $Q_D^{\theta_2} \subseteq
    Q_D^{\theta_1}. $ \hfill \qed


\autoref{tab:example_stream_database} shows a database
of three streams from a cluster monitoring
application~\cite{reiss2012towards}. Each event has three
attributes, denoting the ID, status, and priority of a job
that is executed on a compute cluster. Now, consider the
query $q = \langle
 (x_0,\text{schedule},\text{low}),$
 $(\_,\text{schedule},\_),$ $(x_0,\text{evict},\text{low})
\rangle$. It matches streams $S_2$ and $S_3$, as, e.g.,
the three query terms of $q$ can be mapped to
events $e_{32}$, $e_{34}$, and $e_{35}$ of
$S_3$ (the mapped events are highlighted in blue). Note that
these are the only possible mappings for this scenario, since
the order of the query needs to be preserved.
Yet, the query does not match $S_1$, as the third query term
cannot be mapped to any event ($e_{14}$, highlighted
in red, does not qualify to be mapped).
With a support threshold of
0.6, query $q$ would be in the set of matching queries for
the stream database.

%\begin{wraptable}{r}{6cm}
\begin{table}
\centering
\small
\vspace{-.0em}
\caption{Example stream database.}
\label{tab:example_stream_database}
	\begin{tabular}{l lrrr}
		\toprule
		Stream & Event &  Job Id  &  Status  &  Priority\\
		\midrule
		$S_1$ & $e_{11}$   & \textbf{\color{cyan}\underline{1}} &
		\textbf{\color{cyan}\underline{schedule}} &
		\textbf{\color{cyan}\underline{low}}  \\
	 & $e_{12}$   & 1 & kill & low  \\
	 & $e_{13}$   & 1 & \textbf{\color{cyan}\underline{schedule}} &
	 high  \\
	 & $e_{14}$   & \textbf{\color{cyan}\underline{1}} &
	 \textbf{\color{red}\dotuline{finish}} &
	 \textbf{\color{red}\dotuline{high}}  \\
		\midrule
		$S_2$ & $e_{21}$   & \textbf{\color{cyan}\underline{2}} &
		\textbf{\color{cyan}\underline{schedule}} &
		\textbf{\color{cyan}\underline{low}}  \\
	 & $e_{22}$   & 3 & \textbf{\color{cyan}\underline{schedule}} &
	 high  \\
	 & $e_{23}$   & \textbf{\color{cyan}\underline{2}} &
	 \textbf{\color{cyan}\underline{evict}} & \textbf{\color{cyan}\underline{low}}
	 \\
		\midrule
		$S_3$ & $e_{31}$   & 4 & schedule & high  \\
	 & $e_{32}$   & \textbf{\color{cyan}\underline{5}} &
	 \textbf{\color{cyan}\underline{schedule}} &
	 \textbf{\color{cyan}\underline{low}}  \\
	 & $e_{33}$   & 4 & finish & high  \\
	 & $e_{34}$   & 6 & \textbf{\color{cyan}\underline{schedule}} &
	 low  \\
	 & $e_{35}$   & \textbf{\color{cyan}\underline{5}} &
	 \textbf{\color{cyan}\underline{evict}} &
	 \textbf{\color{cyan}\underline{low}}\\
		\bottomrule
\end{tabular}
%\vspace{-1em}
\end{table}




% \subsection{Query Discovery}
% \label{sec:discovery}

% Given a stream database, the problem of event query discovery relates to
% the identification of event queries that match to a sufficient number of streams to
% fulfill the support threshold.
% % , i.e., for
% % which there exists at least a single match for each stream.
% %are interested in queries for which there exists at least a single match of
% %the query for each stream.
% However, there may exist queries that fulfill the support threshold,
% but which are comparable in the sense that one of them
% is stricter than another one. Formally, a query $q$ is defined as stricter
% than a query $q'$, denoted by $q\prec q'$, if
% % $M(S,q)\subset M(S,q')$ for \emph{any} event stream $S$.
% %$\forall$ Stream $S:$ if $S \models q \Rightarrow S \models q', \exists \
% %S$ s.t. $S \models q$, but $S \not\models q$.
% (i) for any possible stream $S$ (not necessarily contained in a given stream
% database), $S \models q$ implies $S
% \models q'$, and (ii) there exists a stream $S$, such that $S \models q'$,
% but $S \not\models q$.

% In event query discovery, we are only interested in the strictest queries
% that fulfill the support threshold. The reasoning being that these queries
% denote the most concise characterization of the patterns of events
% that
% indicate a situation of interest. For a stream database $D$, we therefore
% consider a notion of descriptiveness~\cite{icdt2022}: A query $q$ is descriptive for $D$,
% if it is supported by $D$ and
% there does not exist a query $q'$ that is stricter, $q'\prec q$, and that is
% also supported by the stream database $D$.
% Based thereon, we formulate the problem of event query discovery:



% \begin{problem}[Event Query Discovery]
% \label{problem}
% Given a stream database $D=\{S_1,\ldots, S_d\}$ and a support threshold $\theta$,
% the problem of event query discovery is to construct the set of queries $Q_D^{\theta}$,
% such that:
% \begin{itemize}[nosep,left=1em]
%   \item $Q_D^{\theta}$ is correct: each $q\in Q_D^{\theta}$ creates at least a single match for
%   a sufficient number streams thus fulfilling the support threshold $\theta$;
%   , i.e., for all $q\in Q_D^{\theta}$ and $S\in M_+(q,D)$ it holds that $S \models
%   q$ and supp$(q,D) = \frac{|M_+(q,D)|}{|D|} \geq \theta$;
%   \item ${Q_D^{\theta}}$ is descriptive: only the strictest queries are considered,
%   i.e., for all $q,q'\in Q_D^{\theta}$ it holds that neither $q\prec q'$ nor $q'\prec q$;
%   \item ${Q_D^{\theta}}$ is complete: if a query $q$ is both correct and descriptive,
%   then it holds that $q\in Q_D^{\theta}$, i.e., ${Q_D^{\theta}Q}$ contains all queries
%   that are both correct and descriptive.
%   \end{itemize}
% \end{problem}








\section{Analysis of the Problem Context}
\label{sec:updates}
% Updating a stream database involves the following two operations: deleting a stream and inserting a stream. The following section describes in detail for each operation which
% information needs to be updated in the new set of descriptive queries.

In this section, we study the problem of event
query discovery under updates, focusing especially on
the relation of the sets of matching queries before
and after an update. We first characterize the problem
context. Let $D_0=\{S_1,\ldots, S_d\}$ be a stream
database,
$\theta_0$ be a support threshold, and
$Q_{D_0}^{\theta_0}$ be a matching query set. For this
setting, we distinguish two possible updates, as
follows, which may also be combined:
\begin{itemize}[nosep,left=.2em .. 1.5em]
  \item The stream database $D_0$ is updated resulting
  in a new database $D_1$. Here, we consider the
  addition of $k$ new
  streams, i.e., $D_1 = D_0 \cup \{S_{d+1}, \dots,
  S_{d+k}\}$, as well as the removal of $k$ existing
  streams, i.e., $D_1 = D_0 \setminus D^-$ with $D^-
  \subseteq D_0$ and $|D^-|=k$.
  \item The support threshold $\theta_0$ is updated
  resulting in a new threshold $\theta_1 \in (0,1]$.
%  \item Both, the stream database $D_0$ and the
%  support threshold $\theta_0$ are updated.
\end{itemize}
These updates potentially have implications for the
set of matching queries, $Q_{D_0}^{\theta_0}$, which,
hence, needs to be updated to $Q_{D_0}^{\theta_1}$,
$Q_{D_1}^{\theta_0}$, or $Q_{D_1}^{\theta_1}$,
respectively. Specifically, we seek to
answer the following questions:
\begin{enumerate}[nosep,left=.2em .. 1.5em]
  \item Which queries are still matching? That is,
  which queries $q\in Q_{D_0}^{\theta_0}$ do still belong to
  $Q_{D_0}^{\theta_1}$,
  $Q_{D_1}^{\theta_0}$, or $Q_{D_1}^{\theta_1}$,
  respectively?
  \item Which new queries are matching? That is,
    which queries $q\notin Q_{D_0}^{\theta_0}$ belong to
    $Q_{D_0}^{\theta_1}$,
    $Q_{D_1}^{\theta_0}$, or $Q_{D_1}^{\theta_1}$,
    respectively?
\end{enumerate}


%\begin{problem}[Event Query Discovery with Updates]
%  \label{problem-updates}
%    Given a stream database $D_0=\{S_1,\ldots,
%S_n\}$, a support threshold $\theta_1$,
%    and a query set $Q_{D_0}^{\theta_0}$, and an
%updated discovery problem.
%    The possible updates include:
%    \begin{itemize}[nosep,left=1em]
%      \item The stream database $D_0$ is updated with
%$k$ new
%      streams $\{S_{d+1}, \dots, S_{d+k}\}$
%      or $k$ existing streams are removed.
%      \item The support threshold $\theta_0$ is
%updated.
%      \item Both the stream database $D_0$ and the
%support threshold $\theta_0$ are updated.
%    \end{itemize}
%    The goal is to update the query set
%$Q_{D_0}^{\theta_0}$ to reflect the changes
%    in the stream database $D_0$ and the support
%threshold $\theta_0$.
%    \begin{enumerate}[nosep,left=1em]
%      \item Which queries in $Q_{D_0}^{\theta_0}$ are
%still correct?
%      \item Which new queries are correct?
%      \item Which information can be derived from the
%previous query discovery?
%      \item How can the query set
%$Q_{D_1}^{\theta_1}$
%      be updated efficiently?
%      \item Is it always feasible to update the query
%      set $Q_{D_1}^{\theta_1}$ incrementally
%      or are we faster by simply running the query
%      discovery on the new setting?
%    \end{enumerate}
%  \end{problem}

%In this section we want to give an answer to the
%question 1-3. Question 4 will
%be answered in \autoref{sec:algos} where we present
%our algorithms.
%Finally, in \autoref{sec:evaluation} we give an
%answer to question 5 comparing
%the runtimes from a baseline event query discovery
%implementation to our
%algorithms that take into account previous
%information.

%\subsection{Relations between Matching Query Sets}
Ideally, the above questions shall be answered
incrementally. That is, we would like to avoid
computing $Q_{D_0}^{\theta_1}$,
$Q_{D_1}^{\theta_0}$, or $Q_{D_1}^{\theta_1}$ from
scratch, but rather derive it directly from
$Q_{D_0}^{\theta_0}$.

Before turning to algorithmic solutions to the above
questions, we study the three relations between the
matching query sets as
illustrated in \autoref{fig:matching_sets}.

\begin{figure}[h]
  %\begin{center}
\centering
\begin{tikzpicture}[scale=2]
    \def\is{1.0} % inner seperation of nodes

    % GROUPS
    \node[inner sep=\is] (Q00) at (0,0) {$Q_{D_0}^{\theta_0}$};
    \node[inner sep=\is] (Q01) at (0,1) {$Q_{D_0}^{\theta_1}$}; % SO+(1,3)
    \node[inner sep=\is] (Q10) at (1,0) {$Q_{D_1}^{\theta_0}$};
    \node[inner sep=\is] (Q11) at (1,1) {$Q_{D_1}^{\theta_1}$};


    % ARROWS
    \draw[<->,thick,myblue] (Q00) -- (Q01)
      node[midway,right=-.65,scale=1.1] {(a)}
      ;
    \draw[<->,thick,myyellow,dotted] (Q00) -- (Q10)
      node[midway,above=-.5,scale=1.1] {(b)}
      ;
    \draw[<->,thick,myblue] (Q10) -- (Q11)
      node[midway,left=-.65,scale=1.1] {(a)}
      ;
    \draw[<->,thick,myyellow, dotted] (Q01) -- (Q11)
      node[midway,below=-.5,scale=1.1] {(b)}
      ;
    \draw[<->,thick,mymagenta, dashdotted] (Q00) -- (Q11)
      node[pos=.75,above =-.75,scale=1.1, rotate=45]
        {(c)}
      ;
  \end{tikzpicture}
  \caption{Illustration of matching query sets and the
  relations between them.}
  \label{fig:matching_sets}
%\end{center}
\end{figure}

In this paper, we focus on the case that $D_1$ is derived from
$D_0$
by the addition of $k$ streams, i.e., $D_1=D_0 \cup
\{S_{d+1}, \dots, S_{d+k}\}$, and that the support
threshold is strengthened, i.e., $0 < \theta_0 \leq
\theta_1 \leq 1$. The reason being that the properties
discussed below hold analogously for the mirrored
updates, i.e., the deletion of existing streams and a
relaxed support threshold.
In the remainder, we first study the relation of the
sets of matching queries
when only the support threshold is updated (a), and
when only the stream database is updated (b).
Afterwards, we turn to the case that both are updated
at the same time (c).


%\textbf{Remark.} Here we only proof the cases that
%streams are added to the
%stream database and the support threshold is
%increased. But the
%porperties hold analogously for the reverted
%scenarios, that is deleting
%streams from the stream database and decreasing the
%support threshold.

%Let $D_0=[S_1,\dots,S_n]$ be a stream database and
%let $D_1$ be an updated stream database
%with $D_1=D_0 \cup \{S_{d+1}, \dots, S_{d+k}\}$.
%    Let the support thresholds $\theta_0, \theta_1$
%be fixed and $0 < \theta_0 \leq \theta_1 \leq 1$ and
%    let $s(|D_0|,\theta_0),s(|D_1|,\theta_1)$ be the
%stream thresholds for $D_0$ and $D_1$, respectively.
%    Let $Q_{D_0}^{\theta_0}$ be the set of matching
%queries for $D_0$ and $\theta_0$ and
%$Q_{D_1}^{\theta_1}$
%    be the set of matching queries for $D_1$ with
%$\theta_1$.

%In \autoref{fig:matching_sets} we have illustrated
%the matching query sets for the
%old and the updated discovery problem.

\subsection*{(a): Updating only the Support Threshold}

In case only the support threshold is changed, i.e.,
it is strengthened, we note that there exists a
containment relation between the sets of matching
queries. Based on \autoref{prop:subset}, we conclude
that the set
of matching queries for the increased support
threshold is a
subset of the one for the lower support threshold, i.e.,
$Q_{D_i}^{\theta_1} \subseteq Q_{D_i}^{\theta_0}$, $i\in \{0,1\}$.
%and $Q_{D_1}^{\theta_1} \subseteq Q_{D_1}^{\theta_0}$.

\subsection*{(b): Updating only the Stream Database}
\label{subsubsec:stream_update}
% \todo[inline]{RS: Add proof for lower bound to better explain the visualization at the end of Section 3b.}
Regarding updates of the stream database, i.e., the
addition of streams, we make the
following observation based on the stream
thresholds, i.e., the minimal numbers of streams that
need to match to fulfill the support threshold.
Consider the case that $s(|D_1|,\theta_i) =
s(|D_0|,\theta_i) + k$, $i\in \{0,1\}$, where
$k$ represents the number of added streams.
Intuitively, this means that
\emph{all} newly added streams need to be matched for
a query to remain in the set of matching queries,
under a worst-case assumption. That is, all queries
that \emph{only} match the minimal numbers of streams
to fulfill the support threshold need to match all the
new streams. We note though that, with a support
threshold of $\theta_i=1$, this worst case applies to
all queries in the set of matching queries.
If $s(|D_1|,\theta_i) =
s(|D_0|,\theta_i) + k$, then the set of
matching queries for
the updated stream database can only become smaller,
i.e., it is
a subset of the original set. This observation on a
containment relation between the matching sets of
queries under updates to the stream database is
captured by the following property:
%$Q_{D_1}^{\theta_0}
%\subseteq Q_{D_0}^{\theta_0}$.
%Analogously, we consider the case that
%$s(n+1,\theta_i)
%= s(n,\theta_i)$, $i\in \{0,1\}$, i.e., not even \emph{one} of the newly
%added streams need to be matched for a query to remain in the set of
%matching queries, the
%respective containment relation holds true:
%$Q_{D_1}^{\theta_1} \subseteq Q_{D_0}^{\theta_1}$.
%This can be summarized in the property:
\begin{property}
\label{prop:db_change_inclusion}
$s(|D_1|,\theta_i) = s(|D_0|,\theta_i) + k, i\in\{0,1\} \implies $
$Q_{D_1}^{\theta_i} \subseteq Q_{D_0}^{\theta_i}$.
\end{property}
The proof of this property is directly covered by \autoref{lemma:update},
which captures a more general case as detailed below. For now, we illustrate
what happens if the respective relation between the stream thresholds does
not hold true with an example.

Consider the streams $S_1, S_2, S_3$, as given in
\autoref{tab:example_stream_database}.
Let $D_0=\{S_1,S_2\}$ be the initial stream database and let the support
threshold be $\theta_0 = 0.6$. Then, the stream threshold is
$s(D_0,\theta_0) = 2$, i.e., a query needs to match both streams to be
frequent. Now, let $D_1=\{S_1,S_2,S_3\}$ be the stream database after adding
one stream, $S_3$, while leaving the support threshold unchanged. Then, the
stream threshold is still $s(D_1,\theta_0) = 2$ and, hence,
$s(D_1,\theta_0)\neq s(D_0,\theta_0) +k$ as $k=1$ in our example.

As mentioned above, the query $q = \langle
(x_0,\text{schedule},\text{low}),(\_,\text{schedule},\_),(x_0,\text{evict},\text{low})
\rangle$ matches $S_2$ and $S_3$, but not $S_1$.
This means that $q \in Q_{D_1}^{\theta_1}$ and $q \notin
Q_{D_0}^{\theta_1}$. As such, query $q$ illustrates that the containment
relation between the sets of matching queries does not hold and we conclude
that, if $s(D_1,\theta_0)\neq s(D_0,\theta_0) +k$, then the relation between
$Q_{D_1}^{\theta_1}$ and $Q_{D_0}^{\theta_1}$ depends on the added stream(s).


%What happens if $s_1\neq s_0 +k$?
%
%  Let $D_1=[S_1,S_2]$ be the stream database with the support threshold
%  $\theta_0 = 0.6$ and the stream threshold $s_0 = 2$.
%Furthermore, let $D_2=[S_1,S_2,S_3]$ be the neighboring stream database
%with
%the same support threshold and
%the stream threshold $s_1 = 2$.


%The following query $q = \langle
%(x_0,\text{schedule},\text{low}),(\_,\text{schedule},\_),(x_0,\text{evict},\text{low})
% \rangle$ matches $S_2$ and $S_3$, but not $S_1$.
%This means that $q \in P_{D_2}^{\theta_1}$ and $q \notin
%P_{D_1}^{\theta_1}$.
%
%Conclusion: If $s_1\neq s_0+k$, the relation of the two matching query sets
%depends on the added streams.

%\exampleboxh{Example 1}{
%What happens if $s_1\neq s_0 +k$?\\
%\label{ex1}
%Consider the following three streams
%$S_1, S_2, S_3$:
%
%	\begin{center}
%	\medskip
%	{\footnotesize
%		\begin{tabular}{l lrrr}
%			\toprule
%			Stream & Event &  Job Id  &  Status  &  Priority\\
%			\midrule
%			$S_1$ & $e_{11}$   & \textbf{\color{cyan}1} &
%\textbf{\color{cyan}schedule} & \textbf{\color{cyan}low}  \\
%			 & $e_{12}$   & 1 & kill & low  \\
%			 & $e_{13}$   & 1 & \textbf{\color{cyan}schedule} & high  \\
%			 & $e_{14}$   & \textbf{\color{cyan}1} &
%\textbf{\color{red}finish} & \textbf{\color{red}high}  \\
%			\midrule
%			$S_2$ & $e_{21}$   & \textbf{\color{cyan}2} &
%\textbf{\color{cyan}schedule} & \textbf{\color{cyan}low}  \\
%			 & $e_{22}$   & 3 & \textbf{\color{cyan}schedule} & high  \\
%			 & $e_{23}$   & \textbf{\color{cyan}2} &
%\textbf{\color{cyan}evict} & \textbf{\color{cyan}low} \\
%			\midrule
%			$S_3$ & $e_{31}$   & 4 & schedule & high  \\
%			 & $e_{32}$   & \textbf{\color{cyan}5} &
%\textbf{\color{cyan}schedule} & \textbf{\color{cyan}low}  \\
%			 & $e_{33}$   & 4 & finish & high  \\
%			 & $e_{34}$   & 6 & \textbf{\color{cyan}schedule} & low  \\
%			 & $e_{35}$   & \textbf{\color{cyan}5} &
%\textbf{\color{cyan}evict} & \textbf{\color{cyan}low}\\
%			\bottomrule
%	\end{tabular}}
%	\end{center}
%
%	\medskip
%  Let $D_1=[S_1,S_2]$ be the stream database with the support threshold
%$\theta_0 = 0.6$ and the stream threshold $s_0 = 2$.
%  Furthermore, let $D_2=[S_1,S_2,S_3]$ be the neighboring stream database
%with the same support threshold and
%  the stream threshold $s_1 = 2$.
%
%
% The following query $q = \langle
%
%(x_0,\text{schedule},\text{low}),(\_,\text{schedule},\_),(x_0,\text{evict},\text{low})
% \rangle$ matches $S_2$ and $S_3$, but not $S_1$.
% This means that $q \in P_{D_2}^{\theta_1}$ and $q \notin
%P_{D_1}^{\theta_1}$.
%
%Conclusion: If $s_1\neq s_0+k$, the relation of the two matching query sets
%depends on the added streams.
%}


The above property and our example illustrate that there is a close relation
between the amount of change applied to a stream database, i.e., the number
$k$ of added streams, the stream threshold, and the support threshold. In
fact, the support threshold can be bounded based on the other two
parameters. These bounds, in turn, enable us to derive insights on how
often we can expect \autoref{prop:db_change_inclusion} to be applicable and,
thereby, conclude that the set of matching queries of the updated database
is a subset of the original set of matching queries.

As above, let $D_0=\{S_1,\dots,S_d\}$ and $D_1=D_0 \cup \{S_{d+1}, \dots,
S_{d+k}\}$ be the original and updated databases and let $\theta_0$ be the
support
threshold for both of them. To simplify the presentation, we use the
short-hands $s_0 = s(|D_0|,\theta_0)$ and $s_1=s(|D_1|,\theta_0)$ for the
stream thresholds $D_0$ and $D_1$, respectively.
Then, it holds that:
$$
s_1 = \lceil (d+k) \cdot \theta_0 \rceil \leq \lceil d \cdot \theta_0 \rceil
+ \lceil k \cdot \theta_0 \rceil \leq \lceil d \cdot \theta_0 \rceil + k =
s_0 + k.
$$
Based thereon, and considering the fact
that $s_0 \leq s_1 \leq s_0+k$, we conclude that $s_1 \in \{s_0,\dots, s_0
+k\}$.

Moreover, knowing that $s_0-1 < \theta_0 \cdot d \leq s_0$, we derive the
following bounds for the support threshold based on the stream threshold
$s_0$ of $D_0$:
$$
\frac{s_0-1}{d} < \theta_0 \leq \frac{s_0}{d}.
$$
Analogously, we derive bounds for the support threshold $\theta_0$ based on
the stream threshold $s_1$ of the updated database $D_1$:
$$
\frac{s_1-1}{d+k} < \theta_0 \leq \frac{s_1}{d+k}.
$$
Combining the above statements, we derive bounds for the support threshold
based on the stream threshold of the original database. However,
we need to distinguish two cases, i.e., $s_1 = s_0+k$, which means that all
added streams need to be matched for a query to remain
frequent under the worst-case assumption that it only
matches exactly $s_0$ streams; and $s_1 <
s_0 +k$, which means that only some of the new streams
need to be matched for any query to remain frequent:
\begin{property}
\label{prop:db_change_bounds}
$s_1 < s_0 +k \implies \frac{s_0-1}{d} < \theta_0 \leq \frac{s_0+k-1}{d+k}$
and $s_1 = s_0+ k \implies \frac{s_0+k-1}{d+k} < \theta_0 \leq
\frac{s_0}{d}.$
\end{property}

% \begin{table}[h!]
%   \centering
%   \begin{tabular}{@{}lllllllll@{}}
%   \toprule
%            & \multicolumn{2}{l}{$s_1$:} & \multicolumn{2}{l}{$s_2$:} & \multicolumn{2}{l}{$s_2=s_1$} & \multicolumn{2}{l}{$s_2=s_1+1$} \\ \midrule
%   $s$/$s_2$ & low         & up         & low         & up          & low          & up          & low          & up            \\
%   1        & 0           & 0.25       & 0           & 0.2         & 0            & 0.2         & 0.2          & 0.25          \\
%   2        & 0.25        & 0.5        & 0.2         & 0.4         & 0.25         & 0.4         & 0.4          & 0.5           \\
%   3        & 0.5         & 0.75       & 0.4         & 0.6         & 0.5          & 0.6         & 0.6          & 0.75          \\
%   4        & 0.75        & 1          & 0.6         & 0.8         & 0.75         & 0.8         & 0.8          & 1             \\
%            &             &            & 0.8         & 1           &              &             &              &               \\ \bottomrule
%   \end{tabular}
%   \caption{The lower and upper bounds for the support threshold $\theta_1$ in relation to the stream threshold $s_1$ and $s_2$ for $n=4$.}
%   \label{tab:bounds}
%   \end{table}

% In \autoref{tab:bounds} we have illustrated the bounds for the support threshold $\theta_1$ in relation to
% the stream threshold $s_1$ and $s_2$ for $n=4$.

% The probability that for a given stream database $D_1$ with $|D_1| = n$ and a given
% stream threshold $s_1$ it holds that $s_2 = s_1$ can be calculated as follows:
% $$
% \mathbb{P}[s_2 = s_1 |n,s_1] = \frac{s_1}{n+1} - \frac{s_1-1}{n} = \frac{1}{n} - \frac{s_1}{n(n+1)}.
% $$
% And over all possible stream thresholds $s_1$ we obtain:
% % $$
% % P(s_2 = s |n) = \sum_{s=1}^{n} P(s_2 = s |n,s) = \sum_{s=1}^{n} \left( \frac{1}{n} - \frac{s}{n(n+1)} \right) =
% % 1 - \frac{1}{n(n+1)} \sum_{s=1}^{n} s = 1 - \frac{1}{n(n+1)} \frac{n(n+1)}{2} = \frac{1}{2}.
% % $$

% \begin{align*}
%     \mathbb{P}[s_2 = s_1 |n] = \sum_{s=1}^{n} \mathbb{P}[s_2 = s_1 |n,s_1] &= \sum_{s_1=1}^{n} \left( \frac{1}{n} - \frac{s_1}{n(n+1)} \right) \\
%     &= 1 - \frac{1}{n(n+1)} \sum_{s_1=1}^{n} s_1 \\
%     &= 1 - \frac{1}{n(n+1)} \frac{n(n+1)}{2} = \frac{1}{2}.
% \end{align*}

% In other words, the probability that the stream threshold $s_2$ is equal to the stream threshold $s_1$ is $\frac{1}{2}$.
% This also means that the probability that the stream threshold $s_2$ is equal to the stream threshold $s_1+1$ is $\frac{1}{2}$.

The bounds as stated in \autoref{prop:db_change_bounds} are useful to derive
insights in the applicability of
\autoref{prop:db_change_inclusion}, i.e., they enable
us to understand when $s_1=s_0+k$ holds after $k$ streams
have been added to a stream database. Intuitively, the
bounds on $\theta_0$ characterize the change points
for the stream threshold, depending on the size of a
stream database $d$ and the size of the change $k$.
We visualize the change points as determined by
\autoref{prop:db_change_bounds} in
\autoref{fig:thresholds}. Here, we consider databases
of different sizes, $d\in \{10,50,100,500,1000\}$, to
which changes of size $k\in \{1,2,5\}$ are applied.
The orange ranges characterize the intervals of the
support threshold $\theta_0$ for which it holds that
$s_1 = s_0+k$, i.e., for which we know that the set of
matching queries after the change is a subset of the
original set of matching queries.
The lower bound of the orange range is determined by
the size of the database and the size of the change,
and can be derived from \autoref{prop:db_change_bounds} as follows:
$$
\frac{s_0}{d}> \frac{s_0+k-1}{d+k} \implies \frac{s_0}{d}>\frac{k-1}{k} \implies \theta_0>\frac{k-1}{k}.
$$
From the visualization, we draw two main conclusions:
\begin{itemize}[nosep,left=.2em .. 1.5em]
  \item The higher the support threshold, the more
  likely
  it is that \autoref{prop:db_change_inclusion} is
  applicable, so that the containment relation between
  the sets of matching queries holds true.
  \item The
  larger the size of the change, the less likely the
  property is applicable.
\end{itemize}





%based on the size of the database, make an estimation
%änderung des stream thresholds in abhängigkeit von
%datenbankgröße und support value
%applicability: wann ist
%abhängig vom support und vom verhältnis der größe
%Bounds auf theta erlauben Rückschlüsse für welche
%theta man eine änderung in dem stream threshold hat
%gegeben datenbankgröße, gegeben ein k
%Daraus folgt die Anwendbarkeit von Property 2
\begin{figure}[h]
  \centering
  \includegraphics[width=.75\textwidth]{img/support_distribution_all.pdf}
  \caption{Visualization of the change points for the
  stream threshold in relation to the support
  threshold, for different scenarios of database sizes
  and change sizes.}
  \label{fig:thresholds}
\end{figure}



% But the probability distribution of the support threshold $\theta_1$ in relation to the stream threshold $s_1$ is not uniform.


% There is a clear trend that the higher the support threshold $\theta_1$
% the more likely it is to fulfill the condition $s_2 = s_1+1$.
% Let_2s consider the case $s_2 = s$.
% Then it holds:
% $$
% \lceil (n+1) \cdot \theta_1 \rceil < \lceil n \cdot \theta_1 \rceil + 1.
% $$
% From this we can deduce that $n \cdot \theta_1 + \theta_1 \leq s$.
% Therefore, $\theta_1 \leq \frac{s}{n+1}$.
% Analogously, we obtain for the case $s_2 = s+1$: $\theta_1 > \frac{s}{n+1}$.



% Then
% \begin{align*}
%     \theta_1 &\leq \frac{\lceil n \cdot \theta_1 \rceil}{n+1} \\
%     (n+1) \cdot \theta_1 &\leq \lceil n \cdot \theta_1 \rceil\\
%     \theta_1 \cdot n + \theta_1 &\leq \lceil n \cdot \theta_1 \rceil \\
%     \theta_1 &\leq \lceil n \cdot \theta_1 \rceil - n \cdot \theta_1.
% \end{align*}
% $$
% \theta_1 \leq \frac{\lceil n \cdot \theta_1 \rceil}{n+1} = \frac{\lceil n \cdot \theta_1 \rceil + n \cdot \theta_1 - n \cdot \theta_1 }{n+1} .
% $$

% \begin{wrapfigure}{r}{0.25\textwidth}
%     \centering
%     \includegraphics[width=0.5\textwidth]{img/support_distribution.pdf}
%     \caption{The support threshold $\theta_1$ in relation to the stream threshold $s$ for $n=4$.}
%     \label{fig:thresholds}
% \end{wrapfigure}



\subsection*{(c): Updating both, the Stream Database
and the Support Threshold}

Finally, we consider the most general case of a
change, where streams are added to the database and
the support threshold is increased. Again, we note
that the change in the stream thresholds that is
induced by the change determines whether the sets of
matching queries are in a containment relation, i.e.,
whether the change leads only to frequent queries
being no longer frequent, but not to additional
frequent queries. We capture this property as follows:
\begin{lemma}
    \label{lemma:update}
    Let $D_0=\{S_1,\dots,S_d\}$ and
    $D_1=D_0 \cup \{S_{d+1}, \dots, S_{d+k}\}$
    be stream databases, and $\theta_0$ and $\theta_1$
    with $\theta_0 \leq \theta_1$ be support
    thresholds for $D_0$ and $D_1$, respectively.
    Then, it holds that $s(|D_1|,\theta_1) \geq
    s(|D_0|,\theta_0) + k \implies Q_{D_1}^{\theta_1}
    \subseteq Q_{D_0}^{\theta_0}$.
\end{lemma}
\textit{Proof.}
Let $q_1$ be a query in $Q_{D_1}^{\theta_1}$. Recall
that $M_+(q_1,D_1)$ is the set of streams in $D_1$
that match $q_1$, $M_+(q_1,D_0)$ is the set of streams
in $D_0$ that match $q_1$, and $M_+(q_1,D_1\backslash
D_0)$ is the set of newly added streams that match
$q_1$. Let $|M_+(q_1,D_1\backslash D_0)|= j$ be the
number of these streams, for which it holds that $j\in
[0,k]$. Then, it suffices to show that $|M_+(q_1,D_0)|
> s(|D_0|,\theta_0)$. As before, we define $s_0 =
s(|D_0|,\theta_0)$ and $s_1 = s(|D_1|,\theta_1)$.
Then, we derive the following bounds for the
cardinality of the aforementioned sets:
$$|M_+(q_1,D_0)| = |M_+(q_1,D_1)| -
|M_+(q_1,D_1\backslash D_0)| \geq s_1 - j \geq s_0 + k
- j \geq s_0.$$
% Then $|M_2_+| \geq s_2 > s_1$.
% If $\{S_{n+1}, \dots, S_{n+k}\}\subseteq M_+(q_1,D_1)$.
% $$|M_+(q_1,D_1)\backslash \{S_{n+1}, \dots, S_{n+k}\}| = s_1 - k
% \geq s_0 + k - k = s_0.$$
% $$|M_2_+\backslash \{S\}| \geq \lceil (n+1) \cdot \theta_1 \rceil - 1
% \geq \lceil n \cdot \theta_1 \rceil = s.$$
% Then $M_+(q_1,D_1)\backslash \{S_{d+1}, \dots, S_{d+k}\} = M_+(q_1,D_0)$ and supp$(q_1,D_0) \geq \frac{|M_+(q_1,D_0)|}{|D_0|} \geq \theta_0$.
% If  $\exists S \in \{S_{d+1}, \dots, S_{d+k}\}$ s.t.
%$S\not\in M_+(q_1,D_1)$:
% If $\{S_{n+1}, \dots, S_{n+k}\}\not\subseteq M_+(q_1,D_1)$.
% Then $M_+(q_1,D_1) \subseteq D_0$.
% $$\frac{|M_+(q_1,D_1)|}{|D_0|} \geq \frac{\lceil (n+k) \cdot \theta_1 \rceil}{n}
% \geq \frac{\lceil n \cdot \theta_0 \rceil}{n} \geq \frac{n\cdot\theta_0}{n} = \theta_0.$$
% Then supp$(q_1,D_0) \geq \theta_0$ and $q_1 \in Q_0$.
Hence, indeed, it holds that $Q_{D_1}^{\theta_1}
\subseteq Q_{D_0}^{\theta_0}$. \hfill  \qed

We note that \autoref{lemma:update} provides a
generalization of \autoref{prop:db_change_inclusion}.
That is, if $\theta_1 = \theta_0$, i.e., streams are
added to the database, but the support threshold
remains unchanged, due to the definition of the stream threshold, $s_0+k$ is an
upper bound for $s_1$. As such, the condition of the implication on the
containment relation becomes $s_1 = s_0+k$.

The above lemma gives a condition, i.e., $s_1\geq
s_0+k$, to conclude on a containment relation for the
sets of matching queries under an update of the stream
database, regardless of the streams that are added.

Considering the general problem context as visualized
in \autoref{fig:matching_sets}, one might assume that
$Q_{D_1}^{\theta_1}\subseteq Q_{D_0}^{\theta_0}
\implies Q_{D_1}^{\theta_0}\subseteq
Q_{D_0}^{\theta_0}$ or $Q_{D_1}^{\theta_1}\subseteq
Q_{D_0}^{\theta_1}$. That is, if the set of matching
queries is only reduced by changes of potentially
both, the stream database and the support threshold,
the containment relation is expected to hold also if
only one of the changes is considered.
However, this is actually not the case, which we
illustrate with a counter example in
\autoref{tab:example}.

\begin{table}[h!]
  \centering
  \begin{tabular}{@{}cccccccc@{}}
  \toprule
  $d$ & $d+1$ & $\theta_0$ & $\theta_1$ &
  $s(d,\theta_0)$ & $s(d+1,\theta_0)$ &
  $s(d,\theta_1)$ & $s(d+1,\theta_1)$ \\ \midrule
  99  & 100   & 0.8        & 0.9        & 80              & 80                & 90              & 90                \\ \bottomrule
  \end{tabular}%
  \caption{Stream thresholds under an example change
  in the stream database and the support threshold.}
  \label{tab:example}
  \end{table}

Here, we consider the setting of a stream database of
size $d=99$, for which one stream is added to obtain a
database of size $d+1=100$. Moreover, we consider a
change of the support threshold from $\theta_0=0.8$ to
$\theta_1=0.9$. For this setting, it holds that
$s(d+1,\theta_0) = s(d,\theta_0)$ as well as
$s(d+1,\theta_1) = s(d,\theta_1)$, i.e., the stream threshold does not change, if
though a stream was added to the database. Hence, without information on this
stream, we cannot
conclude that $Q_{D_1}^{\theta_0} \subseteq
Q_{D_0}^{\theta_0}$ or
$Q_{D_1}^{\theta_1} \subseteq Q_{D_0}^{\theta_1}$.
However, the change of both, the stream database and the support threshold,
actually leads to a change in the stream threshold, i.e., $s(d+1,\theta_1) \geq
s(d,\theta_0)$. Based thereon, using \autoref{lemma:update}, we can conclude on
the containment relation for the sets of matching queries, $Q_{D_1}^{\theta_1}
\subseteq Q_{D_0}^{\theta_0}$.

% We can now take a look at the probability of the different cases for the stream thresholds $s_2$ and $s_1$,
% namely $s_2 = s_1$ and $s_2 \geq s_1+1$.

% With $s_1 = \lceil n \cdot \theta_1 \rceil$ and $s_2 = \lceil (n+1) \cdot \theta_2 \rceil$ we know the boundaries for
% $\theta_1$ and $\theta_2$ in relation to the stream thresholds $s_1$ and $s_2$:

% \begin{align*}
%      \frac{s_1-1}{n} &< \theta_1 \leq \frac{s_1}{n} \\
%     \frac{s_2}{n+1} < &\theta_2 \leq \frac{s_2}{n+1}.
% \end{align*}

% Subtracting the two inequalities we obtain:
% $$\frac{s_2-1}{n+1} - \frac{s_1}{n} < \theta_2 - \theta_1 < \frac{s_2}{n+1} - \frac{s_1-1}{n}.$$

% Now we can look at the two cases:\\
% $s_2 = s_1$:
% \begin{align}
%      \frac{s_1-1}{n} - \frac{s_1}{n} &< \theta_2 - \theta_1 < \frac{s_1}{n+1} - \frac{s_1}{n+1} \nonumber\\
%      \implies -\frac{1}{n+1} - \frac{s_1}{n(n+1)} &< \theta_2 - \theta_1 < \frac{1}{n} - \frac{s_1}{n(n+1)} \nonumber \\
%       \implies 0 &< \theta_2 - \theta_1 < \frac{1}{n} - \frac{s_1}{n(n+1)}. \label{eq:case1}
%     \end{align}

%   $s_2 \geq s_1+1: s_2=s_1+k+1, k\in[0,1,\dots,n-1]$:
%   \begin{align}
%     \frac{s_1+k}{n+1} - \frac{s_1}{n} &< \theta_2 - \theta_1 < \frac{s_2}{n+1} - \frac{s_1-1}{n} \nonumber \\
%     \implies \frac{k}{n+1} - \frac{s_1}{n(n+1)} &< \theta_2 - \theta_1 < \frac{1}{n} + \frac{k+1}{n+1} - \frac{s_1}{n(n+1)} \label{eq:case2} \\
%   k=0:  \implies 0 &< \theta_2 - \theta_1 < \frac{1}{n} + \frac{1}{n+1} - \frac{s_1}{n(n+1)}. \label{eq:case2k0}
%   \end{align}

%   To obtain all possible combinations of $\theta_1$ and $\theta_2$ we have to consider all possible values for $s_1$ and $k$.
%   $$
%   \sum_{s_1=1}^{n} \left( \underbrace{\frac{1}{n} - \frac{s_1}{n(n+1)}}_{\autoref{eq:case1}} +
%    \underbrace{\frac{1}{n} + \frac{1}{n+1} - \frac{s_1}{n(n+1)}}_{\autoref{eq:case2k0}} +
%   \sum_{k=1}^{n+1-s_1} \underbrace{\frac{1}{n} - \frac{s_1}{n(n+1)}}_{\autoref{eq:case2}} \right) = \frac{1}{2} + n - \frac{1}{2} = n.
%   $$

%   For the probabilities of the different cases we obtain:
%   \begin{align*}
%     &\mathbb{P}[s_2 = s_1] = \frac{1}{2n} \\
%     &\mathbb{P}[s_2 \geq s_1+1] = \frac{2n-1}{2n}.
%   \end{align*}


\section{Algorithms for Incremental Discovery of Event Queries}
\label{sec:algos}

In this section, we show how event query discovery is
realized
incrementally, under the changes discussed in the previous section. To this
end, we first summarize a basic discovery procedure that serves as a
baseline (\autoref{sec:discovery}). Then, we discuss how
to update
the set of matching queries (i) once streams are added
and/or the support threshold is increased (\autoref{sec:insert_increment}),
and (ii) once
existing streams are deleted and/or the support threshold
is decreased (\autoref{sec:delete_decrease}).

\subsection{Basic Algorithm for Event Query Discovery}
\label{sec:discovery}

As a basic algorithm for event query discovery, we consider an approach that
searches through the space of candidate queries to collect all matching
queries. To this end, it follows an iterative generate-and-test scheme. For
each query that matches the stream database, stricter child queries are
generated through the addition of query terms, or the insertion of variables
and attribute values for placeholders in existing
query terms. The child
queries are, in turn, evaluated over the stream
database. If they match, they become part of the
result set and their child
queries are considered further.

By starting the approach with an empty query, the approach iteratively
explores the search space of candidate queries. It terminates once all child
queries of all matching queries have been assessed.

The essence of this approach is captured by the function \textsc{Discovery},
formalized in \autoref{alg-discovery}.
It takes as input a stream database $D$; a support threshold $\theta$;
the frequent stream alphabet $\Gamma(D,\theta)$;
a set of queries $Q$ that are frequent in $D$; a parent dictionary
$P$ that is modelled as a
function $P: \mathcal{Q} \to \mathcal{Q}$, mapping child queries to their
parent; and a stack $O$ of queries that are frequent in $D$ and still shall
be explored by a depth-first search through the
space of candidate queries.

\begin{algorithm}[t]
	\small
	\caption{\textsc{Discovery}}
	\label{alg-discovery}
	\KwIn{\, \ Stream database $D$,
		%event schema $\mathcal{A}$, sequence of query var. $\mathcal{X}^<$,
		support threshold $\theta$,
		frequent stream alphabet $\Gamma(D,\theta)$,
		 \\ \hspace{3em} \,
		set of queries $Q$, parent dictionary $P$, query stack $O$
		% \\ \hspace{3em} \,
		% Boolean flag $b\in \{\True,\False \}$ whether to insert variables
		.}
	\KwOut{Updated set of queries $Q$, updated parent dicitionary $P$.}
	\BlankLine

	\tcc{\scriptsize Explore queries that are pushed to the stack}
	%${O \leftarrow}$ empty stack; $O.\textsc{push}(q)$\;
	\While{${O}$ not empty}
	{
		${q \leftarrow O}.\textsc{pop}()$\;
		\If{\normalfont{\textsc{Match}}$({q,D,\theta}$)} {
			\tcc{\scriptsize For frequent queries,
				explore their children}
			${Q \leftarrow Q \cup \left\{q\right\} }$\;
			${P \leftarrow}$
			\textsc{ChildQueries}($\mathit{q,\Gamma(D,\theta),P}$)\;
			${C} \leftarrow  \{q'\in \dom(P) \mid P(q')=q\}$\;
			\If{$ C \neq \emptyset$}{
				\lForEach{$q'\in C$}{
					$O.\textsc{push}(q')$}
			}
		}
	}
	{\Return{$\mathit{Q,P}$}}
\end{algorithm}

Function \textsc{Discovery} then explores the queries on the given stack.
Each of them is evaluated against the stream database using the Boolean
function \textsc{Match}, which returns true if query $q$ is frequent in the
stream database $D$ as defined by the support threshold $\theta$; otherwise,
it returns false. If the query is frequent, it is added to the result set
$Q$ and its child queries are generated using function \textsc{ChildQueries}
and recorded in the aforementioned parent dictionary $P$. If such child
queries have indeed been generated, i.e., set $C$ is not empty, these
queries are pushed to the stack for further exploration.
When the stack is empty, the function returns the set of matching
queries as well as the updated parent dictionary.

Moreover, we summarize the functionality of \textsc{ChildQueries} to
generate child queries of a query $q$ as follows. It generates child queries
by the insertion of a \emph{single} symbol, i.e., a variable or an attribute
value from the frequent stream alphabet, for a placeholder in one of the
existing
query terms. In addition, it
also generates further child queries by appending a new query term to the
query that is empty (i.e., comprises placeholders) except for one position
that is filled with a symbol (variable or attribute value). In general,
these insertions consider all possible attribute values, as well as a set of
possible variables (for which the size is bounded by half of the length of
the shortest stream in the database).

While we use function \textsc{Discovery} also in our incremental approaches
for event query discovery, a basic discovery algorithm simply instantiates
by adding the empty query to the query stack and creating its entry in the
parent dictionary. We formalize this instantiation in
\autoref{alg-baseline}. Given a stream database $D$ and support threshold
$\theta$, it initializes the result set $Q$, derives the frequent stream
alphabet $\Gamma(D,\theta)$, constructs the empty query
$\langle\varepsilon\rangle$ and pushes it to the stack, before starting the
exploration with function \textsc{Discovery}.

\begin{algorithm}[H]
	\small
	\caption{\textsc{Basic Discovery Algorithm} }
	\label{alg-baseline}
	\KwIn{\, \ Stream database $D$, support threshold $\theta$.}
	\KwOut{Set of matching queries $Q$, updated parent dictionary $P$.}
	\BlankLine
	$\mathit{Q \leftarrow \emptyset}$\;
	$\Gamma(D,\theta) = \{a \in
	\Gamma \mid |\{S \in D \mid \exists \ j\in\{1,\dots,|S|\},
	i\in\{1,\dots,n\}: S[j].A_i = a\}| \geq s(|D|,\theta)\}$\;
%	$\Gamma_D \gets \{a \in \Gamma \mid |M_+(a,D)| \geq \lceil |D| \cdot
%	\theta \rceil\}$\;
	$q \gets\langle\varepsilon\rangle$\;
	$P(q) \gets
	\langle\varepsilon\rangle$\;

	${O \leftarrow}$ empty stack\;
	$O.\textsc{push}(q)$\;
	$ Q,P \leftarrow \textsc{Discovery}(D, \theta, \Gamma(D,\theta), Q, P,
	O)$\;
	{\Return{$\mathit{Q,P}$}}
\end{algorithm}

%The baseline algorithms \autoref{alg-baseline} relies
%on an iterative generation of query candidates,
%which are then matched against a stream database. It follows a
%bottom-up direction in the exploration of candidate queries and step-wise
%generation of stricter queries for a given query $q$. It starts with the
%empty
%query, which will be elaborated further later.
%The child queries of $q$ are those obtained by inserting
%a new variable, by inserting a variable that has been present already, or by
%inserting an attribute value; either way, considering an existing query term
%or a newly added query term.
%\color{blue}
%Variables are inserted in a certain order, which
%is reflected in a total order $<$ over a set of variables $\mathcal{X}$. To
%simplify the notation, we write $\mathcal{X}^<$ for the sequence of
%variables induced by $<$ over $\mathcal{X}$. Also, our
%construction employs at most $s/2$ variables with $s=\min_{S\in D} |S|$
%being the minimal stream length in the stream database. Hence, the size of
%$\mathcal{X}$ and the length of $\mathcal{X}^<$ are bounded.
%
%Query candidate generation, formalized in
%\autoref{alg-bu-df:childquery}, takes as input a query $q$; an event
%schema $\mathcal{A}$; a supported alphabet $\Gamma_D$; a parent dictionary
%$P$ that is modelled as a
%function $P: \mathcal{Q} \to \mathcal{Q}$, mapping child queries to their
%parent; a sequence of variables $\mathcal{X}^<$; and a Boolean flag $b$ to
%control whether variables shall be inserted. The algorithm returns the set
%of child queries for $q$.
%
%\autoref{alg-bu-df:childquery} is divided into three parts, following an
%initialization (\autoref{alg1:init_start} - \ref{alg1:init_end}),
%each capturing one type of insertion. It is designed such that \emph{all}
%child queries are obtained by
%insertion
%of a single symbol (attribute value or variable) are constructed
%\emph{exactly once}. The generation relies on the auxiliary function
%$\textsc{Next}$
%(\autoref{alg1:next_start} - \ref{alg1:next_end}). Given a query, it
%generates query candidates by inserting a given symbol (variable or
%attribute value) in
%any query term following a given position in the query, if the query
%contains a placeholder for that attribute. In addition, it
%generates candidates by appending a new query term to the query that
%is empty except for the given symbol for the given attribute. Based thereon,
%the three parts of the main algorithm include:
%
%\textit{Insertion of a new variable} (\autoref{rule:insertnewvar} -
%\ref{rule:insertnewvar_end}). For each attribute, we consider all query
%terms after the last occurrence of any attribute value ($g$) and after the
%first
%occurrence of the last inserted variable ($f$). Using function
%$\textsc{Next}$, a single new variable $\mathcal{X}^<[v+1]$ is introduced
%by replacing a placeholder or adding a new query term.
%
%\begin{algorithm}[H]
%	\footnotesize
%	\caption{\textsc{ChildQueries}: Query Cand. Generation}
%	\label{alg-bu-df:childquery}
%	\KwIn{\, \ Query $q$, event schema $\mathcal{A}$, supported alphabet
%		$\Gamma_D$, \\ \hspace{3em} \,
%		parent dictionary $P$, sequence of query variables $\mathcal{X}^<$,
%		\\ \hspace{3em} \,
%		Boolean flag $b\in \{\True,\False \}$ whether to insert variables.
%	}
%	%\tcc*{some comment}		\;
%	\KwOut{Updated parent dictionary $P$.}
%	\SetKwFunction{FNext}{\textnormal{\textsc{Next}}}%
%	\BlankLine
%	$\mathit{g \gets}$ position of last inserted attribute value in $q$\;
%	\label{alg1:init_start}
%	$\mathit{f \gets}$ first position of last inserted variable in $q$\;
%	$z \gets \max(g,f)$\;
%	$\mathit{n \leftarrow} |\mathcal{A}|$\;
%	\label{alg1:init_end}
%	\tcc{\scriptsize Shall variables be inserted?}
%	\If{$b= \True$}{
%		$v \gets$ number of distinct query variables in $q$\;
%		\tcc{\scriptsize Insert a new variable
%			$\mathcal{X}^<[v+1]$}
%		\ForEach{$A\in\mathcal{A}$}{
%			\label{rule:insertnewvar}
%			\ForEach{$q'\in
%				\FNext(q,z,A,\mathcal{X}^<[v+1],n)$}{
%				\label{rule:insertnewvar_step2}
%				\ForEach{$q''\in
%					\FNext(q',z+|q'|-|q|,A,\mathcal{X}^<[v+1],n)$}{
%					$P(q'') \leftarrow q$\;
%				}
%			}
%		}
%		\label{rule:insertnewvar_end}
%		\tcc{\scriptsize Insert last inserted variable
%			$\mathcal{X}^<[v]$}
%		$l \gets$ last position of last inserted variable in $q$\;
%		$s \gets$ last inserted symbol in $q$\;
%
%		\If{$s=\mathcal{X}^<[v]$}{
%			\label{rule:insertexvar}
%			$\mathit{A \leftarrow}$ attribute $A\in \mathcal{A}$ with $s\in
%			\dom(A)$\;
%			\lForEach{$q'\in \FNext(q,l,A,\mathcal{X}^<[v],n)$}{
%				$P(q') \leftarrow q$
%			}
%		}
%		\label{rule:insertexvar_end}
%	}
%	\tcc{\scriptsize Insert new supported attribute value}
%	\ForEach{$A\in\mathcal{A}$}{
%		\label{rule:inserttype}
%		\ForEach{$a\in (\Gamma_{D}\cap \dom(A))$}
%		{
%			\lForEach{$q'\in \FNext(q,z,A,a,n)$}{
%				$P(q') \leftarrow q$
%			}
%		}
%	}
%	\label{rule:inserttype_end}
%	\Return{$\mathit{P}$}
%
%
%	\BlankLine
%
%	\tcc{\scriptsize Function to generate the next candidate queries}
%	\FNext{\textnormal{query} $q$, \textnormal{start pos.} $s$,
%		\textnormal{attribute} $A$, \textnormal{symbol}	$y$,
%		\textnormal{number
%			of attributes} $n$}{
%		\label{alg1:next_start}
%
%		$\mathit{Q' \leftarrow \emptyset}$\;
%		\ForEach{$i\in \{s,\ldots, |q|\}$}
%		{
%			\tcc{\scriptsize Generate candidate by replacing placeholder
%				with
%				given variable/attribute value $y$}
%			\If{$q[i].A=\_$}{
%				$q' \leftarrow  q$\;
%				$q'[i].A \leftarrow y$\;
%				$Q' \leftarrow Q' \cup \left\{q'\right\}$\;
%
%			}
%			\tcc{\scriptsize Generate candidate by including new query term
%				that
%				only contains given variable/attribute value $y$}
%			$e' \leftarrow \varepsilon$\;
%			$e'.A \leftarrow y$\;
%			$q'\leftarrow \langle q[1],\dots,
%			q[i], e', q[i+1],\dots, q[|q|]
%			\rangle$\;
%			$Q' \leftarrow Q' \cup \left\{q'\right\}$\;
%		}
%		\Return{$\mathit{Q'}$}
%		\label{alg1:next_end}
%	}
%\end{algorithm}


\subsection{Stream Insertion and/or Increased Support
Threshold}
\label{sec:insert_increment}

Having discussed a basic discovery algorithm, we now
present a first algorithmic solution to incremental
processing under changes. Specifically, we consider the
case that streams are added and/or that the support
threshold is increased.

We summarize the idea behind the algorithm as follows:
From
\autoref{lemma:update}, we know that the set of matching
queries after the update will be a subset of the original
set of matching queries. As such, no new queries need to
be generated; we only need to check whether the
queries in the given set are still
frequent. Moreover, the scope of this verification depends
on the support threshold. If it remains unchanged and only
the database changed, it suffices to evaluate the existing
queries over the newly added streams. Otherwise, if
the threshold is increased as well, the entire updated
stream database needs to be considered.

We formalize this approach in \autoref{alg_gen}. It
takes as input the original stream database $D_0$; the
original support threshold $\theta_0$; the
set of matching queries $Q_0$ obtained for $D_0$ and
$\theta_0$; the parent dictionary $P$ for these
queries; the updated stream database $D_1$; and the
updated support threshold $\theta_1$. It returns the
set of matching queries $Q_1$ for $D_1$ and
$\theta_1$ along with the updated parent dictionary
$P$.

After initializing the result set $Q_1$, the
algorithm determines the scope for the verification
of the existing queries, as the entire updated
stream database $D_1$ or the added streams
$D_1\setminus D_0$. Then, the verification is
operationalized with stack onto
which we push the queries that shall be checked. We start
with the empty query and then go through $Q_0$ by means of
the parent-child-relation between queries, as stored in
the parent dictionary.

%If a queries matches, it is added to the updated query
%set and its children are
%added to the stack of queries to be checked for their
%matching behavior.

% In case it does not match, we we add its parent query to the stack of queries to be checked.
% Then discard the query entry from the parent dictionary $P$.

%Once the stack is empty, the algorithm returns the
%updated query set $Q_1$ and the parent dictionary $P$.


\begin{algorithm}[t]
    \small
    \caption{\textsc{UpdatePlus}}
    \label{alg_gen}
    \KwIn{\, \ Stream database $D_0$, support threshold $\theta_0$,
    set of matching queries $Q_0$,
	 \\ \hspace{3em} \,
    parent dictionary $P$,
    updated stream database $D_1$,
    updated support threshold $\theta_1$.}
    \KwOut{Set of matching queries $Q_1$, parent
    dictionary $P$.}
    \BlankLine
    $\mathit{Q_1 \leftarrow \emptyset}$\;
    \lIf{$\theta_1 > \theta_0$}
    {
        $D \leftarrow D_1$}
    \lElse{$D \leftarrow D_1\backslash D_0$}

%    $\mathit{R \leftarrow \emptyset}$\;
    ${O \leftarrow}$ empty stack\;
    % \lForEach{$q\in Q$}
    %     {$R \leftarrow R \cup \left\{P(q)\right\} $}
    % \lForEach{$q\in Q\backslash R$}{
    % $O.\textsc{push}(q)$}
    $O.\textsc{push}(\langle\varepsilon\rangle)$\;
    \tcc{\scriptsize Explore queries that are pushed to the stack}
    \While{${O}$ not empty}
    {
            ${q \leftarrow O}.\textsc{pop}()$\;
            % \If{$q \not\in Q_1$}{

%            $\mathit{M \leftarrow}$
%\textsc{Match}(${q,D,\theta}$)
%            \;
            % \If{$ M = \False$} {
            \If{\normalfont{\textsc{Match}}(${q,D,\theta}$)}
             {
            \tcc{\scriptsize For frequent queries, add
            the query
            to the result set and
            push their child queries onto the stack}

            ${Q_1 \leftarrow Q_1 \cup \left\{q\right\}
            }$\;
            % $O.\textsc{push}(P(q))$\;
            % $P.\textsc{pop}(q)$ \;
            ${C} \leftarrow  \{q'\in \dom(P) \mid P(q')=q\}$\;
            \lForEach{$q'\in C$}{
            $O.\textsc{push}(q')$}
            }
            % \lElse{
            %     ${Q_1 \leftarrow Q_1 \cup \left\{q\right\} }$}

        % }
    }
    {\Return{$\mathit{Q_1,P}$}}
    \end{algorithm}

\subsection{Stream Deletion and/or Decreased Support Threshold}
\label{sec:delete_decrease}

Next, we turn to incremental event query discovery once
streams are deleted from the database and/or the support
threshold is decreased. Again, our approach relies on
\autoref{lemma:update}, and exploits that all queries in
the existing set of matching queries will
also match the updated stream database.

In general, the idea of handling stream deletion and/or a decreased support
threshold in an incremental manner can be summarized as follows.
We consider the queries, for which the child queries are not frequent.
Intuitively, these queries denote the frontier of frequent queries in the
space of all query candidates. For these queries, we derive the child
queries and add them to the stack of queries to be explored further using
the basic discovery procedure. However, this exploration first considers
solely the attribute values that have been frequent in the original
database. In order to ensure completeness of the set of matching queries, in
a second step, we therefore also check whether the frequent alphabet in the
updated stream database contains new attribute values. If so, we generate
the child queries of all frequent queries found so far, but only incorporate
the new frequent attribute values (and no insertion of variables) in the
respective generation. All queries generated in this way, had not
been considered previously. If one of these queries is frequent, we
iteratively continue the exploration with their child queries, this 
time inserting all frequent attribute values as well as variables.
As such, the algorithm is similar to the basic discovery procedure
introduced in \autoref{sec:discovery}, but shows an important difference: It
only assesses queries that have been not been checked as part of the
previous discovery run over the database before the change.



%For these queries, we construct the child
%calculating their children and adding them to the stack of queries to be
%matched.
%Then the algorithm checks whether the supported alphabet of the new stream
%database is the same
%as the supported alphabet of the old stream database.
%In case, the supported alphabet includes new types, we consider all
%matching queries calculated so far
%and generate new candidate queries by adding the new types to the supported
%alphabet.
%We prevent queries from being generated more than once by first adding to
%the existing queries
%only the newly supported types. All queries generated in this way, had not
%been considered
%previously.
%If a query matches, we keep generating
%(then considering the entire supported alphabet) new candidate queries
%until we reach a non-match.
%% Finally, we disgard all matching queries which are not descriptive to
%%obtain the updated set of descriptive queries.
%The algorithm is similar to the baseline algorithm with the main difference
%that it only checks queries that were not part of the previous event query
%discovery.


We formalize the respective algorithm in \autoref{alg_spec}. The input is
the same as discussed for \autoref{alg_gen}, in terms of the stream
databases, support thresholds, parent dictionary and the set of matching
queries. The algorithm then initializes the frequent alphabets for the
stream databases before and after the change, i.e., $\Gamma(D_0,\theta_0)$
and $\Gamma(D_1,\theta_1)$. Then, we collect all parents of frequent
queries, as a set $R$, to be able to iterate over all queries, for which the
child queries are not frequent, given by $Q_1 \setminus R$. For these
queries, the child queries are generated using function
\textsc{ChildQueries}, considering the frequent alphabet
$\Gamma(D_0,\theta_0)$, and pushed to a query stack. The latter is used in
function \textsc{Discovery}, from \autoref{alg-discovery}, to explore the
space of candidate queries, starting from the queries on the stack.

\begin{algorithm}[t]
    \small
    \caption{\textsc{UpdateMinus}}
    \label{alg_spec}
    \KwIn{\, \ Stream database $D_0$, support threshold $\theta_0$,
    set of matching queries $Q_0$,
    \\ \hspace{3em} \,
    parent dictionary $P$,
    updated stream database $D_1$, updated support
    threshold $\theta_1$.}
    \KwOut{Set of matching queries $Q_1$, updated parent
    dictionary $P$.}
    \BlankLine
    $\mathit{Q_1 \leftarrow Q_0}$\;
    $s_0 \gets s(|D_0|,\theta_0)$\;
    $\Gamma(D_0,\theta_0) \gets \{a \in
    \Gamma \mid |\{S \in D_0 \mid \exists \
    j\in\{1,...,|S|\},
    i\in\{1,...,n\}: S[j].A_i = a\}|\geq s_0\}$\;
    $s_1 \gets s(|D_1|,\theta_1)$\;
    $\Gamma(D_1,\theta_1) \gets \{a \in
    \Gamma \mid |\{S \in D_1 \mid \exists \
    j\in\{1,...,|S|\},
    i\in\{1,...,n\}: S[j].A_i = a\}| \geq s_1\}$\;
%    \Gamma_{D_0} \gets \{a \in \Gamma_0 \mid |M_+(a,D_0)|
%\geq \lceil |D_0| \cdot \theta_0 \rceil\}$\;
%    $\Gamma_{D_1} \gets \{a \in \Gamma_1 \mid
%|M_+(a,D_1)| \geq \lceil |D_1| \cdot \theta_1 \rceil\}$\;
    $\mathit{R \leftarrow \emptyset}$\;
    ${O \leftarrow}$ empty stack\;
    \lForEach{$q\in Q_1$}
        {$R \leftarrow R \cup \left\{P(q)\right\} $}
    \ForEach{$q\in Q_1\backslash R$}{
        ${P \leftarrow}$
        \textsc{ChildQueries}($\mathit{q,\Gamma({D_0},\theta_0),P}$)\;
        ${C} \leftarrow  \{q'\in \dom(P) \mid P(q')=q\}$\;
        \If{$ C \neq \emptyset$}{
        \lForEach{$q'\in C$}{
            $O.\textsc{push}(q')$}
        }
    }
    $ Q_1,P \leftarrow \textsc{Discovery}(D_1,\theta_1,
    \Gamma({D_0},\theta_0), Q_1, P, O)$\;

    \If{$\Gamma({D_1},\theta_1)\backslash
    \Gamma({D_0},\theta_0) \neq \emptyset$}{
        \ForEach{$q\in Q_1$}{
        ${P \leftarrow}$
        \textsc{ChildQueriesNoVars}($\mathit{q,\Gamma({D_1},\theta_1)\backslash
         \Gamma({D_0},\theta_0),P}$)\;
        ${C} \leftarrow  \{q'\in \dom(P) \mid P(q')=q\}$\;
        \If{$ C \neq \emptyset$}{
        \lForEach{$q'\in C$}{
            $O.\textsc{push}(q')$}
        }
    }
    $Q_1,P \leftarrow \textsc{Discovery}(D_1,
    \theta_1, \Gamma({D_1},\theta_1), Q_1, P, O)$\;

    }
    {\Return{$\mathit{Q_1,P}$}}
    \end{algorithm}

As mentioned, we need to consider that an update to the stream
database changed the frequent alphabet, i.e., attribute values became
frequent due to the removal of streams. If so, i.e., if
$\Gamma({D_1},\theta_1)\backslash
\Gamma({D_0},\theta_0)$ is not empty, for each query obtained so far, we
need to consider the child queries constructed
only through the insertion of these new frequent attribute values
(neglecting the insertion of additional variables), which is captured by
function \textsc{ChildQueriesNoVars}. For these queries, further exploration
may be needed, so that they are pushed onto the query stack, which is then
used in function \textsc{Discovery}.

\section{Experimental Evaluation}
\label{sec:evaluation}
% \todo[inline]{R2 on scalability: Note that the scalability is not directly affected by the size of the stream database, $|D|$,
% but rather by the structure of the streams. A small database with complex,
% repetitive patterns can cause combinatorial explosions in candidate queries, leading to runtimes of minutes or hours.
% Conversely, a large database with simple, non-repetitive patterns may be processed in seconds.
% Thus, scalability evaluations based on $|D|$ or the alphabet size $|\Gamma|$ are unreliable, as runtime depends heavily on stream structure.
% One explanation for the difference in the datasets is the size of the queryset.
% If for example, the queryset stays the same after updating the database, then the updated algorithm will terminate quickly
% whereas the baseline algorithm has to rediscover all matching queries. We will add the explanation to the final version of the paper.}

This section presents an evaluation of the algorithms proposed in
\autoref{sec:algos}. We
first introduce our experimental setup (\autoref{sec:exp_setup}), before
comparing our algorithms for incremental event query discovery against a
baseline solution (\autoref{sec:exp_sota}).
% Finally, we evaluate the sensitivity of
% our algorithms with respect to the characteristics of a stream database
% (\autoref{sec:exp_sensitity}), before studying the effectiveness of our
% strategies to guide the application of discovery algorithms
% (\autoref{sec:exp_guidance}).

\subsection{Experimental Setup}
\label{sec:exp_setup}

\sstitle{Datasets}
We use two real-world datasets to evaluate our algorithms, namely
 NASDAQ stock trade events~\cite{eodata}
and the Google cluster traces~\cite{reiss2011google}.
We adopted the methodology for obtaining stream databases from the~\cite{ilminer}.
For each dataset, we define three queries to generate streams, as described below.

The attributes considered in the finance data are:
% The event schema for the finance data comprises three attributes:
stock name, volume, and closing value. The queries to construct streams
capture
two events with stock name `GOOG' (F1); three events with stock names
appearing in the sequence `AAPL', `GOOG', and `MSFT' (F2); and four events
for which the volume remains constant (F3).

Events in the Google cluster traces contain four attributes:
job ID, machine ID, task status, and priority. One query
defines three events running on the same machine with identical job IDs,
with the first and last ones of status `submit', while the
second one has
status `kill' (G1); one query encompasses four events
executing on the same
machine with status `finish' (G2); and one query detects job eviction
resulting from the scheduling of another job on the same
machine (G3).

Using these queries, we generated the stream databases summarized in
\autoref{tab:databases} by evaluating the queries over the
datasets. For each match (of the first at most 1000 matches), we constructed a
stream by including the events preceding the match.
The number of events taken into account for each stream
 is detailed in the column 'stream length'.


\begin{table}[h]
	\caption{Stream databases used in the experiments.}
	\label{tab:databases}
	\small
	% \vspace{-1.2em}
	\centering
	\begin{tabular}{c c c c c }
		\toprule
		\multicolumn{2}{c}{}          & stream
			length& number
			of
			streams\\%& \begin{tabular}[c]{@{}c@{}}support\\
			%threshold\end{tabular} \\
		\midrule
		\multirow{3}{*}{Finance} & F1 & 50                    &
		69-79                                                           \\
		& F2 & 40
		& 990-1000                                                       \\
		& F3 & 25
		& 240-250                                                           \\
		%\midrule
		\multirow{1}{*}{Google}  & G1-G3 & 7              &
		990-1000                                                         \\
		\bottomrule
	\end{tabular}
	%\vspace{-1em}
\end{table}

Note that the runtime is not determined by the size of the stream database,
$|D|$, but rather by the structure of the streams. A small database with intricate and repetitive patterns can trigger
a combinatorial explosion of candidate queries, leading to runtimes of minutes or hours.
Conversely, a large database with simple, non-repetitive patterns may be processed in a matter of seconds.

\sstitle{Baseline}
% \todo[inline]{RS: Add experiments for combined updates.}
We compare the runtime of our algorithms making use of incremental updates against
the baseline approach, in which no prior information is given to the
discovery.


\sstitle{Measures} We primarily measure the efficiency of event query
discovery in terms of the algorithms' runtime. We report averages over five
experimental runs, including error bars (mostly negligible).
%For each stream database, the query discovery was repeated 5 times and the
%average runtime is plotted including the confidence interval error bars.

\sstitle{Environment and Implementation}
We implemented the whole experimental evaluation in Python.
The code and the evaluation routine is publicly available on GitHub\footnote{https://github.com/rebesatt/IncDISCES}.
All experiments were conducted on a server with a E7-4880 CPU @2.5GHz and
1TB of RAM.
\subsection{Baseline Comparison}
\label{sec:exp_sota}

% We first conducted experiments on the datasets described in \autoref{tab:databases}
% by decreasing the space of candidate queries, i.e., by adding streams, by
% increasing the support threshold or both.
We first conducted experiments on the datasets described in \autoref{tab:databases} by reducing
the candidate query space—either by adding streams, increasing the support threshold, or both.

In \autoref{fig:gen_sample}, we plotted the relative time change of the
incremental algorithm
compared to the baseline, when adding up to 10 streams to the different
stream databases.
It can be seen that our incremental algorithm is at least an order of
magnitude faster for all settings.
For the google dataset G1, the improvement is even up to three orders of
magnitude.

In \autoref{fig:gen_supp}, we consider the change of increasing the support
threshold.
The original support threshold is 0.75, which is increased stepwise, up to
1.0.
Here, we see a drop in the relative time change for increasing support
values. This is
especially true for the finance stream databases. One reason for this
observation is the
decreasing search space. For example, the number of queries for F1 and
support threshold
1.0 is four. As consequence, both, the baseline and the incremental
algorithm terminate quickly, in less than a second.

% In \autoref{fig:gen_sampsupp}, we consider the simulatenous change of increasing the support
% threshold and added streams.

In \autoref{fig:gen_sampsupp}, we examine the simultaneous effect
of increasing the support threshold and adding streams.
Like in the previous experiment, the original support threshold is 0.75 and
is increased stepwise up to 1.0. At the same time, for each step, we add
one stream to the stream database.
Similar to \autoref{fig:gen_supp}, we see a drop in the relative time change
for increasing support values, but still the runtime of the incremental algorithm is
always lower than the baseline solution.


% \begin{figure}[t]
% 	\centering
% 	\begin{subfigure}{.63\textwidth}
% 	  \centering
% 	  \includegraphics[width=.82\linewidth]{img/gen_sample.pdf}
% 	  \includegraphics[clip,trim=73em 0em 5.5em 0em,
% 	  width=.14\linewidth]{img/gen_supp.pdf}
% 	  \caption{Adding streams to the stream database}
% 	  \label{fig:gen_sample}
% 	\end{subfigure}\\[1em]%
% 	\begin{subfigure}{.63\textwidth}
% 	  \centering
% 	  \includegraphics[width=\linewidth]{img/gen_supp.pdf}
% 	  \caption{Increasing the support threshold}
% 	  \label{fig:gen_supp}
% 	\end{subfigure}
% 	\begin{subfigure}{.63\textwidth}
% 	  \centering
% 	  \includegraphics[width=\linewidth]{img/gen_sampsupp.pdf}
% 	  \caption{Increasing sample size and the support threshold}
% 	  \label{fig:gen_sampsupp}
% 	\end{subfigure}
% 	\caption{Relative time change of updated algorithm compared to baseline
% 	for different stream databases while decreasing the search space.}
% 	\label{fig:test}
% 	\end{figure}

\begin{figure}[t]
    \centering
    \begin{minipage}[t]{0.48\textwidth} % Left column
        \centering
        \vspace{0pt} % Force alignment at the top
        \begin{subfigure}{\textwidth}
            \centering
            \includegraphics[width=\linewidth]{img/gen_sample.pdf}
            % \includegraphics[clip,trim=73em 0em 5.5em 0em, width=.14\linewidth]{img/gen_supp.pdf}
            \caption{Adding streams to the stream database}
            \label{fig:gen_sample}
        \end{subfigure}\\[1em]
        \begin{subfigure}{\textwidth}
            \centering
            \includegraphics[width=\linewidth]{img/gen_supp.pdf}
            \caption{Increasing the support threshold}
            \label{fig:gen_supp}
        \end{subfigure}
        \begin{subfigure}{\textwidth}
            \centering
            \includegraphics[width=\linewidth]{img/gen_sampsupp.pdf}
            \caption{Increasing sample size and the support threshold}
            \label{fig:gen_sampsupp}
        \end{subfigure}
        \caption{Relative time change of updated algorithm compared to baseline for different stream databases while decreasing the search space.}
        \label{fig:test}
    \end{minipage}
    \hfill
    \begin{minipage}[t]{0.48\textwidth} % Right column
        \centering
        \vspace{0pt} % Force alignment at the top
        \begin{subfigure}{\textwidth}
            \centering
            \includegraphics[width=\linewidth]{img/spec_sample.pdf}
            % \includegraphics[clip,trim=76em 0em 5.5em 0em, width=.14\linewidth]{img/spec_supp.pdf}
            \caption{Deleting streams to the stream database}
            \label{fig:spec_samp}
        \end{subfigure}\\[1em]
        \begin{subfigure}{\textwidth}
            \centering
            \includegraphics[width=\linewidth]{img/spec_supp.pdf}
            \caption{Decreasing support threshold}
            \label{fig:spec_supp}
        \end{subfigure}
        \begin{subfigure}{\textwidth}
            \centering
            \includegraphics[width=\linewidth]{img/spec_sampsupp.pdf}
            \caption{Decreasing sample size and support threshold}
            \label{fig:spec_sampsupp}
        \end{subfigure}
        \caption{Relative time change of updated algorithm compared to baseline for different stream databases while increasing the search space.}
        \label{fig:test2}
    \end{minipage}
\end{figure}



% Next, we evaluated the behaviour of the algorithms when increasing the space
% of candidate queries that need to be searched, by deleting streams from the
% database and/or by decreasing the support threshold.
Next, we evaluated the behaviour of the algorithms when increasing the space
of candidate queries that need to be searched, by deleting streams from the
database and/or by decreasing the support threshold.

In \autoref{fig:spec_samp}, for each dataset, the
incremental
algorithm outperforms the baseline solution. For the google stream database
G2, the
incremental algorithm is twice as fast independently of the number of
streams deleted
from the original database. The highest benefit is observed for the google
dataset G3,  where the updated algorithm is up to 50 times faster.

In \autoref{fig:spec_supp}, we see that for decreasing support threshold for
all datasets, the two algorithms need roughly the same
runtime to discover the queries.
This is because the original setting has a small set of matching queries. As
a consequence, the benefit of the incremental discovery algorithms is also
small, which yields a similar runtime as the baseline
solution.

Similarly, in the experiment shown in \autoref{fig:spec_sampsupp},
where both the support threshold and the number of streams are decreased,
the results follow the same trend. The small set of matching queries in the original
setting limits the advantage of the incremental discovery algorithms, leading to comparable
runtimes with the baseline approach.


% \begin{figure}[t]
% 	\centering
% 	\begin{subfigure}{.63\textwidth}
% 		\centering
% 		\includegraphics[width=.82\linewidth]{img/spec_sample.pdf}
% 		\includegraphics[clip,trim=76em 0em 5.5em 0em,
% width=.14\linewidth]{img/spec_supp.pdf}
% 		\caption{Deleting streams to the stream database}
% 		\label{fig:spec_samp}
% 	\end{subfigure}\\[1em]%
% 	\begin{subfigure}{.63\textwidth}
% 		\centering
% 		\includegraphics[width=\linewidth]{img/spec_supp.pdf}
% 		\caption{Decreasing support threshold}
% 		\label{fig:spec_supp}
% 	\end{subfigure}
% 	\begin{subfigure}{.63\textwidth}
% 		\centering
% 		\includegraphics[width=\linewidth]{img/spec_sampsupp.pdf}
% 		\caption{Decreasing sample size and support threshold}
% 		\label{fig:spec_sampsupp}
% 	\end{subfigure}
% 	\caption{Relative time change of updated algorithm compared to baseline
% 	for different stream databases while increasing the search space.}
% 	\label{fig:test2}
% 	\end{figure}
\section{Related Work}
\label{sec:related-work}
%\todo[inline]{We acknowledge that techniques for
%sequential rule mining are related.
%Yet, there is a fundamental difference in the adopted
%query model.
%In the revised version, we will add a more
%comprehensive discussion of these techniques,
%clarifying their differences and the unique
%challenges in the context of CEP.}
Our work relates to existing approaches for automated query discovery, as
well as incremental algorithms for sequential pattern mining and frequent
itemset mining.

\textbf{Automated Query Discovery}
So far research in the area of automated query discovery did not take into account
updates in the database or the support threshold.
We adopted the query model used in ~\cite{DBLP:conf/btw/Kleest-Meissner23}. But
the proposed algorithm does not find all matching queries but only discovers
one matching query for each run.
The other approaches \cite{ilminer,icep} further lack in the query
expressiveness. While \cite{icep} cannot discover queries with
multiple occurrences of an attribute value, \cite{ilminer} lacks
the ability to find query terms which only contain variables.

\textbf{Incremental Sequential Pattern Mining}
A related field in which incremental algorithms have been proposed
is sequential pattern mining. Here, the input is a database of strings
rather than a stream of events which can contain multiple attributes.
Furthermore, the discovery is limited to type level without the
ability to discover variables or patterns.
In \cite{zheng2002,zheng2003}, the authors propose two algorithms, one for
adding
and one for deleting data to the original database. The algorithms make use of
negative border sequences which are those sequences whose sub-sequences
are frequent to recalculate the results.
\cite{mallick2013} presents a survey that summarizes the incremental
approaches in this
area. There the highlighted algorithms are based on common frequent
sequence mining algorithms.





\textbf{Incremental Frequent Itemset Mining}
In frequent itemset mining sets of items are discovered that occur frequently
together in a database. This setting differs from our query discovery since
the events in a discovered query are order sensitive whereas item sets can occur
in arbitrary order.
High-utility pattern mining, a sub-field of frequent itemset mining,
discovers itemsets based on a user-defined utility threshold.
The incremental algorithm introduced in \cite{yun2015} proposes a tree
structure to update the item sets.

\textbf{Sequential Rule Mining}
Sequential rule mining \cite{fournier2011rulegrowth,erminer} is a combination of sequential pattern mining and
frequent itemset mining. The goal is to find rules that describe a
relationship between itemsets in a sequence. The rules are of the form
$X \rightarrow Y$ where $X$ and $Y$ are itemsets. The rules are
evaluated based on metrics like support and
confidence, and my also be found
incrementally~\cite{ierminer}.
However, one limitation of such rules is that they are
bounded to
one implication whereas our queries can contain multiple events.
Also, the rules are limited in their expressiveness as
the itemset captures only type-level information, not
allowing for variables.



% Theory paper recommended by Sarah
% \cite{idris2020}

\section{Conclusions}
\label{sec:conclusions}

In this paper, we focused on the problem of event query discovery in the presence of
changes. That is, once all frequent queries have been discovered for a given stream
database and a given support threshold, we studied how this set can be updated
incrementally if the database changes and/or the support threshold changes, in order to
avoid to recompute the respective queries from scratch. We first provided a theoretical
analysis of the problem context. Here, we could establish conditions under which the set
of matching queries after the change is a subset of the respective set before the
change or vice versa. This information is valuable as it highlights when it is sufficient to verify
whether the existing queries are still frequent, but there is no need for the generation
of additional queries. Based on these insights, we also presented algorithmic solutions
that adapt the set of queries incrementally. Our experiments using real-world data
demonstrate that our approach to incremental query
discovery is significantly faster than a baseline that recomputes the queries, reducing
runtime by up to three orders of magnitude. In future work, we intend to
increase the efficiency further,
e.g., through the design of data structures that capture further details on where in the
streams the queries match as well as through the use of statistics about the stream
database (e.g., on the distribution of attribute values over streams) and we plan to extend
our evaluations by adding alternative performance metrics, such as memory usage.


% \begin{acks}
\section*{Acknowledgments}
 This work was funded by the German Research
Foundation (DFG), CRC 1404: "FONDA: Foundations of
Workflows for Large-Scale Scientific Data Analysis".
%  Additional funding was provided by [...] and [...].
% We also thank [...] for contributing [...].
% \end{acks}

%%% Angabe der .bib-Datei (ohne Endung) / State .bib file (im Falle der
%%%Nutzung von BibTeX)
\bibliography{references}
% \printbibliography % im Falle der Nutzung von biblatex
\end{document}
